% Created 2026-06-23 Tue 17:53
% Intended LaTeX compiler: pdflatex
\documentclass[11pt,letterpaper]{article}
\usepackage[utf8]{inputenc}
\usepackage[T1]{fontenc}
\usepackage{graphicx}
\usepackage{longtable}
\usepackage{wrapfig}
\usepackage{rotating}
\usepackage[normalem]{ulem}
\usepackage{amsmath}
\usepackage{amssymb}
\usepackage{capt-of}
\usepackage{hyperref}
\usepackage{amsmath}
\usepackage{amssymb}
\usepackage{geometry}
\geometry{left=1.2in, right=1.2in, top=1.0in, bottom=1.0in}
\usepackage{fancyhdr}
\usepackage{booktabs}
\usepackage{array}
\usepackage{setspace}
\usepackage{parskip}
\usepackage{microtype}
\PassOptionsToPackage{hidelinks}{hyperref}
\newcommand{\Kappa}{\mathcal{K}}
\setlength{\parindent}{0pt}
\setlength{\parskip}{6pt}
% --- Page style: no header on title/TOC, header from page 3 onward ---
\fancypagestyle{plain}{\fancyhf{}\fancyfoot[C]{\thepage}\renewcommand{\headrulewidth}{0pt}}
\fancypagestyle{body}{%
\fancyhf{}%
\fancyhead[C]{\small Version 0.5 $\cdot$ DOI: [pending] $\cdot$ CC BY 4.0}%
\fancyfoot[C]{\thepage}%
\renewcommand{\headrulewidth}{0.4pt}}
\author{Bradley K. DePuy, Independent Researcher}
\date{2026}
\title{The Medium: A Foundational Theory of Reality}
\hypersetup{
 pdfauthor={Bradley K. DePuy, Independent Researcher},
 pdftitle={The Medium: A Foundational Theory of Reality},
 pdfkeywords={},
 pdfsubject={},
 pdfcreator={Emacs 29.3 (Org mode 9.6.15)}, 
 pdflang={English}}
\begin{document}


%% ---- TITLE PAGE (page 1, no header) ----
\pagestyle{plain}
\begin{titlepage}
\vspace*{4cm}
\begin{center}
  {\LARGE\bfseries The Medium}\\[0.6em]
  {\large A Foundational Theory of Reality}\\[2.5em]
  {\normalsize Bradley K. DePuy, Independent Researcher}\\[0.5em]
  {\normalsize ORCID: \mbox{\texttt{0009-0001-0328-2299}}}\\[0.5em]
  {\normalsize 2026}
\end{center}
\end{titlepage}

%% ---- Body begins, header active ----
\pagestyle{body}

\section{Introduction}
\label{sec:orgaf0a4fc}
\subsection{The Core Problem}
\label{sec:org04eb0ff}


Physics inherited a set of primitives. Object, location, motion, sequential time, the detached external observer. These were abstracted from human-scale perceptual experience — the cave dweller's sky, formalized. Every subsequent framework inherited them without examination. Thermodynamics, electromagnetism, relativity, quantum mechanics, string theory, loop quantum gravity. Each built on the same foundation.

The incompatibilities between frameworks — GR against QM, the measurement problem, the problem of time — are symptoms of foundational mismatch. Not technical failures awaiting mathematical fixes.

\subsection{The Language Problem}
\label{sec:org353be6a}


Quantum mechanics borrowed ordinary language and bent each term. Observer no longer means a conscious mind. Measurement no longer means careful observation. Collapse may not be a physical event at all. The borrowed terms carry associations the new meanings cannot support. False familiarity misleads and gatekeeps critique.

The Medium requires its own vocabulary. Three terms replace the inherited ones:

\begin{itemize}
\item \emph{Occurrence}: a \(\iota\) event — the unit of substrate process
\item \emph{Registration}: retention of a structural consequence of interaction
\item \emph{Boundary}: the interface at which two processes alter each other's state
\end{itemize}

Chosen to mean exactly what they say. Nothing more.

\subsection{A Recognition and a Challenge}
\label{sec:org9d896a7}

Modern science is one of the great achievements of the human mind. Quantum mechanics predicts the behavior of matter at the smallest accessible scales with a precision unmatched in the history of empirical inquiry. General relativity describes the geometry of spacetime and the behavior of gravity across cosmological distances with equal exactness. The Standard Model catalogs the fundamental particles and their interactions with extraordinary completeness. These are not approximations or guesses. They are hard-won structures --- the accumulated result of centuries of observation, mathematical invention, and experimental testing.

This document does not dispute any of it.

What it asks instead is the question: \emph{What alternative foundation explains the findings of modern science?} The equations of physics are precise maps. The challenge is to identify the territory they map --- not the territory of observable phenomena, which physics has surveyed magnificently, but the foundation that gives rise to a universe in which observation is possible at all.

There is a limitation that physics has not confronted fully. Anyone who has ever done physics is an earth-bound observer, even with the advent of space-borne sensors. That observer arrived late --- billions of years after the assumed beginning of the universe --- in the midst of an apparently infinite universe. That observer, equipped with sensory tools required for survival, began to ask the questions. Using a dimensional coordinate system (meters, seconds, joules), physics constructed systems defined to understand what is seen from that observation point. Constants, dimensions, and equations carry the fingerprints of that position.

\section{Abstract}
\label{sec:org5a29088}

The Medium is a foundational framework prior to the limits of human observation. It cannot be tested or demonstrated directly. It is offered as a starting point upon which others may build.


The sole primitive is \(\iota\): the Medium's act of self-distinction. The Medium distinguishes within itself. \(\iota\) is that act. Each \(\iota\) event produces character --- the difference the distinction makes. Character is not a property \(\iota\) holds. It is what distinguishing is. Without difference there is no distinction. \(\iota\) is not point-like. It has multiple relational modes --- qualitatively distinct aspects of the \(\iota\) event itself, not binary opposition. \(\iota\) has intrinsic cohesive attraction: not because it is complementary to another \(\iota\) event, but because attraction is what \(\iota\) is. From intrinsic attraction, self-organization follows without instruction. The Medium is the endless occurrence of \(\iota\), unconstrained. No boundary, structure, or assumed constraint exists at the substrate level. Space, time, dimensionality, and causality are not conditions imposed upon \(\iota\) --- they emerge from \(\iota\)-relations. The projection boundary marks the threshold at which the substrate meets the observer's dimensional framework. In the Medium's unit system, one \(\iota\) event = 1. What SI physics calls \(\hbar\) is the SI measurement of that unit --- not a fundamental constant, but a unit conversion artifact locating the projection boundary on the human-scale SI grid.

The central organizing claim of v0.5: processes occur independently. No ordering. No coordination required. Yet existence may impact existence. Sequence is what this looks like from inside an embedded process. The substrate field \(\phi\) records distinction characters across the Medium.

From distinction characters, a three-level algebraic hierarchy emerges by successive combination. A scalar account derives \(\kappa = \sum \iota\) and \(\Kappa = \sum \kappa\), with \(\Kappa\) projecting into spacetime as \(\hat{H}\), the Hamiltonian of standard physics --- not assumed but derived from substrate structure. That scalar account is provisional: the multifaceted \(\iota\) model requires richer combination rules --- more than addition, because the modes differ in kind. The Hamiltonian endpoint survives revision.

Valid combinations are governed by a logical adjacency constraint: two \(\iota\) events of identical distinction character produce no new information and therefore cannot occupy adjacent positions in a combination. This is not a physical rule imposed on \(\iota\) but a logical necessity --- identical adjacent distinction characters produce nothing new. The combinatorial space of valid configurations is sharply restricted by this constraint. A minimum threshold of combinatorial complexity must be reached before any stable configuration can form --- below this threshold, combinations approach but do not achieve coherence. At the view layer, these sub-threshold combinations appear as quantum vacuum fluctuations. The quantum vacuum has measurable effects. Whether my model gives the right numbers, I don't yet know.

Stable configurations emerge where the relational web reaches a stability extremum. These are the relational analogs of bound states. Mass, charge, and spin are not independent intrinsic properties --- they are three projection signatures of a single underlying relational structure. Mass arises at two levels: at the substrate level as a stable stochastic imperfection in the \(\iota\) packing; at the view layer, configuration scale corresponds to \(N\) as a first approximation. The specific mass values and the hierarchy between generations are not yet derived. Charge corresponds to the net distinction character of a configuration. Spin describes the relational closure of a configuration's encounter pattern within the web. It is not rotation.

The Medium distinguishes three categories of emergent property. Properties emergent from the substrate arise directly from \(\iota\)-combinations: mass, charge, spin, the quantum vacuum, the Hamiltonian, spacetime, and dimensionality. Properties emergent at the view layer arise from interactions between stable configurations in spacetime: the fundamental forces, fields, and the full complexity of observable reality. This is the domain of existing physics --- The Medium makes no claims there. A third category remains unassigned: neutrino oscillation and dark matter have not yet been placed with confidence. Dark energy is a projection artifact of misread redshift, not a substance. Gravity has a candidate direction: mass is a stable stochastic imperfection in the \(\iota\) packing; gravity is the strain field propagating through the medium as it attempts to correct toward ground state. Infinite reach without a carrier particle is the expected consequence; formal development is open. Color charge is a labeling convention, not a substrate property. The physics is the stability pattern, not the label. The GR/QM conflict is a category error. Three layers: substrate granularity, relational-web averaging, view-layer smoothness. Quantizing gravity fails because gravity is already the averaged-out result.

The relational web is the probability landscape of accumulated stochastic \(\iota\) interactions. The wave function is the view layer's representation of that landscape. The Born rule is how the view layer reads probability density. Wave function collapse is not a physical event.

The four major puzzles of standard cosmology --- the entropy problem, the arrow of time, cosmological redshift, and the singular origin of the Big Bang --- are addressed in Part III. In each case the puzzle dissolves rather than being solved, because it was generated by a category error that The Medium does not commit: treating view-layer artifacts as substrate-level facts.

Since \(\iota\)-arrangements are below the current observational floor, direct testing is not possible with existing instruments. The framework predicts that particle properties treated as intrinsic by standard physics are relational outcomes of substrate combinatorics. New experimental approaches designed to probe transition signatures rather than particle states directly are identified as the path toward indirect verification.

The central open problems are offered as challenges to the theoretical physics community: deriving the \(\iota\) relational mode structure and its combination rules; rebuilding the path from substrate to Hamiltonian via the mode structure; deriving the specific configurations corresponding to observed particles; establishing which relational patterns produce which spin values under projection; completing the formal account of the gravity strain field; and characterizing dark matter candidates within the substrate.

\emph{Changes from v0.4:}
\begin{itemize}
\item Reorganized around the independence primitive as the central organizing claim
\item Revised \(\iota\): occurrence without attribute; attribute produced at encounter, not carried by \(\iota\)
\item Encounter established as the primitive event from which relational character and particle properties derive
\item Added Part IV: What The Medium Needs to Prove
\item Added new vocabulary layer (occurrence, registration, boundary) to replace inherited QM terminology
\item Added Observer and Observed section dissolving the observer/observed asymmetry at the substrate level
\item Expanded time section with the real-for-the-observer vs. necessary-for-the-universe distinction
\item Added The Core Problem and The Language Problem as Introduction subsections
\item Revised \(\iota\): not point-like; multiple relational modes, not binary opposition; intrinsic cohesive attraction as a property of \(\iota\) itself
\item Mass identified as a stable stochastic imperfection in the \(\iota\) configuration; gravity as the strain field propagating through the medium toward ground-state restoration
\item Relational web reframed as a probability landscape --- stochastic all the way down, statistical consequence not active process
\item GR/QM conflict dissolved across the three layers: substrate granularity, relational-web averaging, view-layer smoothness
\item Wave function and collapse dissolved via the probability landscape account
\end{itemize}


\section{PART I: THE SUBSTRATE}
\label{sec:orgad00cb6}

The substrate description is complete in itself. It contains no dimensions, no observers, no physics. It is what it is.

\subsection{The Independence Primitive}
\label{sec:orgd8b802d}


Processes occur independently. No ordering. No coordination required. Yet existence may impact existence.

Sequence is what independence looks like from inside an embedded process. Experienced as before and after. The experience is real. The sequence is not the substrate.

Identity is produced at encounter. Each process occurs on its own terms. The relationships are what the processes are. Remove the web and there is no process. Remove the process and there is no web.


Three prior frameworks reach toward this ground and stop short.

String theory places a string at the foundation --- a one-dimensional entity carrying vibration modes intrinsically. Each frequency corresponds to a different particle. Spacetime is already there before the string. It is the container. Properties are in the entity, not the encounter.

Loop quantum gravity makes spacetime emergent from a spin network --- nodes and edges that are atoms of geometry. That is a genuine step. But the network structure itself is assumed, not derived. LQG does not explain why that relational structure exists. It starts one level too late.

Planck bits --- Wheeler's \emph{It from Bit}, Fredkin, Wolfram --- place a discrete bit at the foundation carrying 0 or 1 intrinsically. The binary distinction is assumed before any encounter. The question of which value is answered before the question is formed.

\(\iota\) is none of these. All three prior frameworks place entities first and derive relationships. The Medium places the act of distinction first. Not an entity. Not a web requiring prior connection. The Medium distinguishes within itself. \(\iota\) is that act. Character is what the act produces.

Planck bits come closest --- both are discrete, dimensionless, pre-spatial. The gap is not small. A bit carries its state intrinsically. \(\iota\) produces character as the act of distinction. That is the difference between assuming distinction and deriving it.

Loop quantum gravity comes closest structurally --- both make spacetime emergent from relational structure. But LQG presupposes the relational structure. The Medium derives it from the act of self-distinction.

\subsection{The Foundational Claim}
\label{sec:org80ba654}

The framework rests on a single foundational claim:

\begin{quote}
The Medium always is and always was and always will be. It has no beginning and no end. Intrinsic distinction is an eternal property of the Medium --- not an event that occurred but a fundamental condition. The structure we observe is not the result of something that happened. It is the expression of the Medium always happening.
\end{quote}

This claim does three things at once.

\textbf{It eliminates the need for an initial starting point.} The question of what initiated the universe only arises with the assumption of time. If the Medium always is and always was, there is no before and no starting point. The trigger needed for a beginning dissolves by identifying it as a consequence of a false assumption.

\textbf{It eliminates time as a primitive.} The Medium exists outside of time. What we call time is not a property of the Medium. It is downstream --- a consequence of repeated manifestation, not a fundamental feature of reality.

\textbf{It re-frames differentiation.} The emergence of structure is not a process that started. It is a condition that always is. The ground state and the differentiated state coexist within the Medium eternally, because intrinsic distinction --- \(\iota\) --- is eternal.

\subsection{The Ground State}
\label{sec:orgfa9dac0}

The Medium at its ground state is a fully autonomous stochastic system. It requires no inputs, no external cause, no temporal framework, and no governing distribution. Its randomness is not randomness within a framework --- it is randomness without framework. No privileged outcome. No preferred state.

This is not chaos. Chaos retains structure. The ground state of the Medium precedes structure entirely. It is meta-randomness --- randomness that is random about its own nature.

In a medium of infinite extent with no preferred states, every possible configuration has nonzero probability. This is not intuitive but is mathematically unavoidable. The universe is not something that happened against the odds. It is a necessary expression of a medium that has no mechanism for suppressing any possibility.

\subsection{Axioms}
\label{sec:org416b6ce}

\subsubsection{Axiom I --- The Eternal Medium}
\label{sec:org9cca792}

The Medium always is and always was. It has no beginning, no end, and no external cause.

\begin{equation}
\nexists\, t_0 : M \text{ begins at } t_0 \qquad \nexists\, t_1 : M \text{ ends at } t_1
\end{equation}

\emph{Status: Axiomatic.}

\subsubsection{Axiom II --- Intrinsic Distinction}
\label{sec:org289ca4a}


The Medium possesses one property: intrinsic distinction, denoted \(\iota\). It is not an operator, not a function, not a process. It is what the Medium is. The Medium distinguishes within itself. \(\iota\) is that act of self-distinction. Each \(\iota\) event produces character as the difference the distinction makes. Without difference there is no distinction. Character is not a property \(\iota\) holds --- it is what distinguishing is.

\emph{Status: Axiomatic.}

\subsection{Primitives}
\label{sec:org3c61cb7}

\subsubsection{Primitive I --- The Field}
\label{sec:org447e321}


\(\iota\) is the Medium's act of self-distinction. Each \(\iota\) event produces character --- the difference the distinction makes. That character is not a property \(\iota\) holds before the act. It is what the act is. +1 and -1 are the two aspects of one distinction. It has no spatial extent, no duration, no location, no mass, no charge, no spin. It is dimensionless.

Each \(\iota\) event occupies a relational node --- a position defined entirely by its co-manifestation relations. Not by spatial coordinates. Spatial position is what a relational node looks like under projection.

The substrate field \(\phi\) records the distinction character produced by each \(\iota\) event. M here is the set of relational nodes, not a spatial manifold.

\begin{equation}
\phi : M \to \{+1, -1\}
\end{equation}

\begin{center}
A record of distinction characters. The complete description of the Medium's state.
\end{center}

At the view layer, one \(\iota\) event = 1 in the Medium's unit system. The SI physics community locates this threshold at the Planck scale --- the scale at which quantum and gravitational effects are simultaneously relevant. That location is an empirical correspondence, not a definition.

\(\iota\) events are not driven. No prior state determines which relational node the Medium next distinguishes within itself. They simply occur.

\emph{Status: Revised in v0.5.}

\subsubsection{Primitive II --- Distinction and Character}
\label{sec:org1ceaf52}


Each \(\iota\) event is an act of self-distinction within the Medium. The character it produces --- +1 or -1 --- is the difference the distinction makes. Not determined in advance. Not driven by prior state. Whether the Medium next distinguishes at a given relational node is not governed. It simply does.

\begin{equation}
\phi(\iota) \in \{+1, -1\} \text{ (character of distinction, no preference, no guarantee)}
\end{equation}

The substrate imposes no preference on which character is produced. Without differential distinction character, \(\iota\) events produce only undifferentiated noise. No structure, no stability, no material. Differential character is the logical necessity for material existence.

A note on the binary axiom. The choice of \(\{+1, -1\}\) is defended on logical grounds only: one act of distinction creates one boundary; one boundary has exactly two sides. That argument does not appeal to quantum mechanics or observed discreteness at the Planck scale. Those observations are downstream --- consistency checks, not justification. To justify binary \(\iota\) by reference to quantum discreteness would embed the observable into the foundation the framework is attempting to derive. The axiom stands or falls on the logical claim alone. Whether that claim is sufficient --- whether one act of distinction necessarily produces exactly two values rather than three or a continuum --- is an open question carried in Unresolved Questions.

Which character a \(\iota\) event produces is not determined before the event. The question does not exist until the distinction occurs. This is not a matter of unknown information. The answer does not exist yet. There is no fact about the outcome until the act of distinction makes one.

\begin{center}
Shared state. No distance. No space. No ordering.
\end{center}

This is prior to locality. Any spatial relationship between \(x\) and \(y\) is constructed downstream, after structure formation establishes the conditions in which distance becomes meaningful.

\emph{Status: Revised in v0.5.}

\begin{quote}
Not a beginning and an end. A breathing.
\end{quote}

Because \(\iota\) is an eternal property of the Medium, this arc does not occur once. It occurs wherever and whenever \(\iota\) manifests. Encounter patterns form, achieve stability, dissolve --- across an infinite relational web that does not notice.

\subsection{The Algebraic Hierarchy}
\label{sec:org906e3dc}

The substrate must connect to existing physics at the view layer. The endpoint is the Hamiltonian \(\hat{H}\) --- the operator that governs how configurations evolve in projected time. That endpoint is required.

The path from substrate to Hamiltonian passes through the relational web. How \(\iota\) events combine, what combination rules govern the relational modes, and how stable configurations emerge from those rules --- that is relational web work, not substrate work. The full account is developed in Part II.

\emph{Status: Endpoint required. Path developed in Part II.}

\subsection{The Physical Character of \(\iota\)}
\label{sec:orgb901a5f}

The \(\iota\) event is not point-like. A point has no internal structure. A point-like \(\iota\) could vary in magnitude or type but not in kind.

\(\iota\) has multiple relational modes. Not the binary \(\{+1,-1\}\) of the field --- those are the distinction characters produced at encounter. The relational modes are qualitatively distinct aspects of the \(\iota\) event itself. Not binary opposition. Not a single axis. Multiple distinct modes --- as quark color charge is not magnitude but three qualitatively distinct states that combine under specific rules. The number of modes is not assumed. It is constrained by what the framework must produce. This connects directly to the open question in Unresolved Questions: whether \(\{+1,-1\}\) is the full description of \(\iota\)'s internal structure, or whether the relational modes require richer characterization.

\(\iota\) has intrinsic cohesive attraction. Not because it is complementary to another \(\iota\) event. Not because it is missing something an opposite supplies. Attraction is what \(\iota\) is. A single \(\iota\) event in isolation has this property. No second event is required to activate it. The attraction is mode-dependent --- different relational modes produce different kinds of attraction, not merely different magnitudes.

From intrinsic attraction, self-organization follows without instruction. \(\iota\) events draw toward each other. The arrangement falls into the lowest-energy configuration the relational modes permit. The geometry emerges from the constraint --- not assumed, not designed. In three-dimensional projection this corresponds to the class of energy-minimizing space-filling structures: the arrangements nature defaults to wherever energy seeks equilibrium under packing constraints.

No packing process is perfect. At the \(\iota\) scale, stochastic fluctuations are inevitable. Some regions configure slightly differently than others. A relational mode misaligned. An arrangement slightly off. These imperfections are guaranteed by probability alone.

The imperfections persist and propagate. A local misalignment strains neighboring \(\iota\) events. They adjust. The adjustment strains their neighbors. The disturbance propagates outward through the medium.

\begin{quote}
Mass is a stable stochastic imperfection in the \(\iota\) configuration.
\end{quote}

A persistent local deviation from the ground-state arrangement that the surrounding medium continuously attempts to correct.

\begin{quote}
Gravity is the strain field propagating through the medium as it attempts to correct toward ground state.
\end{quote}

Not a force transmitted across a void. Not a signal. The medium's continuous attempt at restoration, propagating through the \(\iota\) network itself. Infinite reach follows without residue. Every \(\iota\) event participates in the ground state. Every stable imperfection propagates strain outward through all \(\iota\) events. No carrier particle required. The participation is constitutive.

\emph{Status: Candidate direction developed 2026-05-23. Requires formal development.}

\subsection{The Substrate --- Complete}
\label{sec:org4d99238}

The substrate description is now complete. It contains:

\begin{itemize}
\item One substrate: the Medium
\item One organizing claim: the independence primitive --- processes occur independently; existence may impact existence
\item One property: \(\iota\) --- the Medium's act of self-distinction, producing character as the difference the distinction makes
\item One field: \(\phi \in \{+1, -1\}\) --- a record of distinction characters
\item One relationship: co-manifestation \(x \sim y\)
\item One derivation target: the path from \(\iota\) interactions to the Hamiltonian, developed in Part II
\item One arc: maximum undifferentiation \(\to\) structure \(\to\) maximum undifferentiation
\end{itemize}

\begin{quote}
Nothing else belongs here. Everything else is downstream.
\end{quote}


\section{PART II: THE RELATIONAL WEB}
\label{sec:org7014147}

The relational web is what emerges from \(\iota\) interactions at scale. It is not a designed structure. It is a statistical consequence. No rules are imposed from outside. The structure is what intrinsic attraction does when \(\iota\) events are many, simultaneous, and stochastic.

The relational web sits between the substrate and the view layer. It is structured but pre-dimensional. It has no spatial coordinates, no temporal ordering, no dimensional quantities. It presents to the view layer only what it can stably support.

\subsection{The Probability Landscape}
\label{sec:org3bca6be}

The relational web is stochastic. Not deterministic. Not programmed. Stochastic all the way down --- from the substrate through the relational web. Probability is the actual character of both levels, not a limitation of knowledge.

The relational web is the probability landscape of accumulated stochastic \(\iota\) interactions. The weight of simultaneous \(\iota\) events creates regions of higher and lower likelihood. No mechanism required. Statistics operating at scale --- as temperature emerges from molecular motion without any molecule possessing temperature.

The relational web feeds back into the substrate without active intervention. Accumulated pattern shapes what comes next --- not because anything decides, but because the worn path makes that route more likely.

What the view layer reads as a particle, a mass, a field is a region of high probability density in that landscape --- stable enough, persistent enough, dense enough to resolve into something nameable. Stable stochastic imperfections persist. Transient fluctuations cancel in the aggregate and wash out.

\subsection{The Combination Rules}
\label{sec:org703abe4}

The relational web governs how \(\iota\) events combine. The combination rules are not imposed from outside. They emerge from the \(\iota\)'s relational modes and intrinsic attraction.

A prior account derived a scalar hierarchy: \(\kappa = \sum \iota\) and \(\Kappa = \sum \kappa\), projecting to \(\hat{H} = \mathrm{proj}(\Kappa)\). That account assumed \(\iota\) produces scalar values combining by addition. The multifaceted \(\iota\) model makes scalar summation insufficient. Relational modes are qualitatively distinct. The combination rules governing them may not be addition. The scalar hierarchy is a first approximation that captured one dimension of \(\iota\)'s structure. It requires rebuilding from the mode structure upward.

The endpoint is not in question. Whatever the combination rules are, they must project to the Hamiltonian \(\hat{H}\) at the view layer --- the view layer's account of how configurations evolve in projected time. The Hamiltonian is a required derivation target, not an assumption.

The derivation of the full combination rules from the \(\iota\) mode structure is carried in Unresolved Questions.

\emph{Status: Endpoint required. Path under revision.}

\subsection{The Adjacency Constraint}
\label{sec:orgbedebf2}

Two \(\iota\) events of identical distinction character cannot be directly co-manifested. This is a logical necessity --- identical adjacent distinction characters produce no new distinction and therefore no new structure. The constraint is not a physical rule imposed on \(\iota\). It is what distinction is.

The binary formulation applies directly to the \(\{+1,-1\}\) distinction characters. Whether the multifaceted \(\iota\) model requires the constraint to be reformulated in terms of relational modes --- and what "identical" means when \(\iota\) events have internal mode structure --- is an open question carried in Unresolved Questions.

What the constraint produces in the relational web:
\begin{itemize}
\item Sharply restricts which configurations are possible
\item Sets a minimum complexity threshold below which stable configurations cannot form
\item Sub-threshold combinations appear at the view layer as quantum vacuum fluctuations
\end{itemize}

\emph{Status: Binary formulation current. Mode-structure reformulation open.}

\subsection{The GR/QM Dissolution}
\label{sec:org46253b8}

General relativity and quantum mechanics have resisted unification for a century. Every attempt to quantize gravity fails. Every attempt to embed quantum mechanics in curved spacetime breaks down at the Planck scale. The conflict is treated as a technical problem awaiting the right mathematical framework.

It is not a technical problem. It is a category error.

GR and QM are not competing descriptions of the same thing. They are accurate descriptions of different layers.

The substrate is genuinely granular. Discrete \(\iota\) events. Discrete interactions. No determinism. Probability is the actual character of the substrate, not a limitation of knowledge about it.

The relational web is the probability landscape of those stochastic interactions. Still granular in origin. Still stochastic. But accumulated over massive numbers of simultaneous \(\iota\) events, producing regions of density.

The view layer is smooth. The averaging process produces continuity. GR's spacetime manifold is real --- at the view layer. It is what the relational web's probability landscape looks like to an embedded observer constructing a dimensional coordinate account.

Quantum mechanics is substrate granularity leaking through to the view layer. The wave function is the view layer's representation of the probability landscape the relational web produces. The Born rule --- probability proportional to amplitude squared --- is how the view layer reads probability density. Wave function collapse is not a physical event. It is the view layer resolving one reading from a probability landscape that was never definite.

General relativity is the smooth aggregate. Spacetime curvature is the view-layer account of large-scale probability density gradients in the relational web.

\begin{quote}
Quantizing gravity fails because gravity is already the averaged-out result. You cannot granularize something that only exists because granularity was averaged away.
\end{quote}

The mathematical failure is the signal. Both frameworks are correct. They are correct at different levels. The conflict between them is the conflict of levels mistaken for a conflict of theories.

\emph{Status: Candidate dissolution developed 2026-05-23. Requires formal development.}

\subsection{The Relational Web --- Complete}
\label{sec:orgcf7f4db}

The relational web description is complete for this version. It contains:

\begin{itemize}
\item One character: stochastic probability landscape --- statistical consequence of \(\iota\) interactions, not imposed structure
\item One constraint: adjacency --- identical distinction characters cannot be directly co-manifested
\item One stability condition: persistent high-density regions in the probability landscape
\item One derivation target: combination rules from \(\iota\) mode structure → Hamiltonian
\item One consequence: GR and QM describe different levels of the same landscape
\end{itemize}

\begin{quote}
The relational web is what intrinsic attraction does at scale. Nothing more is required.
\end{quote}

\section{PART III: THE VIEW LAYER}
\label{sec:org604f760}

Everything in Part III is downstream of the relational web. Dimensions, constants, observers, and all of physics live here --- valid, precise, and complete as descriptions of projected structure. Not descriptions of the substrate or relational web itself.

\subsection{The Three-Level Structure}
\label{sec:orgdb83f4b}

The framework has three levels, not two.

\textbf{Level 1 --- The Substrate.} \(\iota\) events. The Medium distinguishing within itself. Raw distinction --- +1 or $-$1 --- occurring without pattern, without structure, without dimensionality. The \(\phi\) field records the distinction characters produced. The adjacency constraint is the one logical rule that applies here. Nothing else belongs at this level.

\textbf{Level 2 --- The Relational Web.} Registered \(\iota\) distinction characters organized by the combination rules into co-manifestation relations. Stable configurations at stability extrema under the combination rules governing the \(\iota\) relational modes. This level is structured but pre-dimensional.

The relational web is stochastic. Not deterministic. Not programmed. Stochastic all the way down --- from the substrate through the relational web. Probability is the actual character of both levels, not a limitation of knowledge.

The relational web is the probability landscape itself. The accumulated weight of stochastic \(\iota\) interactions creates regions of higher and lower likelihood. No mechanism required. Statistics operating at scale. Like temperature emerging from molecular motion: no single molecule has temperature, but the aggregate of molecular motion is temperature. The relational web does to \(\iota\) interactions what temperature does to molecular motion --- it collapses granular activity into a readable signal.

The relational web has structural consequence without active intervention. Accumulated pattern shapes what comes next --- not because anything decides, but because the worn path makes that route more likely. The aggregate feeds back into the substrate purely by existing.

What the view layer reads as a particle, a mass, a field is a region of high probability density in that landscape --- stable enough, persistent enough, dense enough to resolve into something nameable. Stochastic imperfections that are stable become mass. Transient fluctuations cancel in the aggregate and wash out.

\textbf{Level 3 --- The View Layer.} The relational web as experienced by an embedded observer. The observer is itself a stable \(\iota\) configuration complex enough to sustain a systematic account of what the relational web presents. That account takes the form of a dimensional coordinate system: length, time, mass, charge. The stable configurations the relational web presents are read as particles. The relational geometry is read as spacetime. The constants of physics --- \(c\), \(\iota_{\mathrm{proj}}\), \(G\) --- are the seams where the observer's dimensional framework meets the pre-dimensional relational web.

The view layer is where physics operates. It is valid, precise, and complete as a description of what the relational web presents to an observer. It is not a description of the relational web itself, and not a description of the substrate.

This three-level structure resolves several questions the two-level account left open. Why do particles have the properties they do? Because stable \(\iota\) configurations at stability extrema have those properties at the relational web level --- the observer reads them as mass, charge, and spin at the view layer. Why does high-energy collision not reliably reveal substrate structure? Because collisions operate at the view layer --- smashing view-layer particles together and reading the debris. The debris consists of view-layer forced states resolving back to what the relational web stably supports. The relational web level is not directly accessible to view-layer observation.

\emph{Status: Established in v0.5.}

\subsection{The Observer's Position}
\label{sec:org1b00a49}

Before describing the projections, it is necessary to name what the observer is and what it is not.

An observer is a self-sustaining structure in the binary field --- a stable pattern that persists across transitions, embedded in the differentiated region it constitutes. Every observer who has ever done physics is such a structure. Every human observer is an earth-bound variant: carbon-based, temperature-sensitive, confined to a narrow spatial and temporal window, equipped with sensory apparatus tuned to a small fraction of the available signal, and dependent on cognitive architecture shaped by survival pressures that predate science by orders of magnitude.

These are not complaints. They are structural facts that bear on what physics can and cannot see from its current position.

The observer arrives late. The substrate has been operating --- in whatever sense ``operating'' applies to an eternal, timeless ground --- without observers for the overwhelming majority of the arc. The observer arrives embedded in a mature, highly differentiated region of the binary field. It has no access to the substrate prior to structure formation. It has no access to regions beyond the propagation horizon of its own local structure. Its coordinate system --- meters, seconds, joules, kelvins --- was built to make \emph{its own scale} tractable. That coordinate system is extraordinarily powerful within its domain. It is not a description of the substrate. It is a description of the projection.

The constants of physics --- \(h\), \(c\), \(G\), \(k\) --- are the seams where the observer's dimensional framework meets the substrate's dimensionless operations. They are not brute inputs. They are artifacts of the position from which the measurement is taken.

A different observer, embedded in a different region of the binary field at a different stage of the arc, with a different sensory range and a different coordinate system, would derive different constants and call them fundamental. They would be equally right and equally downstream.

\subsection{Physics as Projection}
\label{sec:orged4991b}

The Medium is pre-dimensional. The binary field \(\phi\) carries no units. The transition index \(t\) is a dimensionless label. \(c\) does not exist at the substrate level. The Independence Primitive establishes that processes occur independently with no shared temporal framework. A rate requires counting transitions against time. No substrate-level time exists. Therefore no substrate-level rate exists, and \(c\) as a transition rate limit is a category error at the substrate level. \(c\) is view-layer emergent --- the ratio of projected distance to projected time, fixed by the structure of projection itself.

\begin{quote}
Dimensions do not exist in the substrate. They are constructed by observers.
\end{quote}

To describe the relationships between structures like itself, the observer constructs a coordinate system --- a framework of units and dimensions that makes the relationships tractable. That coordinate system is not a discovery about the substrate. It is a construction of the observer.

\begin{quote}
Every dimensional quantity in physics --- every constant, every observable, every equation --- is a projection of a substrate relationship onto the observer's dimensional coordinate system. Physics is the complete and consistent description of those projections. It is a perfect map of the shadow. It is not a description of the substrate that casts it.
\end{quote}

\subsection{The Projection Boundary and Natural Units}
\label{sec:org8928bfc}

The Planck scale marks the projection boundary: the threshold at which one \(\iota\) event becomes visible to the observer's dimensional framework.

In the Medium's unit system, one \(\iota\) event is 1. Not 1 of something. Just 1 --- the irreducible unit of occurrence. No smaller unit exists, because no smaller event exists.

SI units cannot be the foundation. They were built from human-scale standards --- the meter, the kilogram, the second --- and carry the fingerprints of the observer's position. To write substrate quantities in SI is to import view-layer categories into a level that has none. Planck units do not escape this. They are just SI constants rearranged, still built from quantities the observer measured. They rescale the observer's framework. They do not replace it.

At the substrate there are no dimensions at all --- substrate quantities are pure counts. The dimension the observer calls \emph{action} is not in the substrate. It is assigned at projection, when the observer wraps one projected event in a dimensional framework.

What physics calls Planck's constant \(\hbar\) is the observer's measurement of that one projected event. It is not a fundamental constant. In the Medium's own units the projected unit is simply 1. Its constancy is not a law of physics: every \(\iota\) event is the same event, so the observer always measures the same thing and gets the same number. It is just 1 equaling 1. The Medium's unit of action is:

\begin{equation}
[\mathcal{A}] \equiv \iota_{\mathrm{proj}} = 1
\end{equation}

From this unit and the structure of projection, the minimum units of length, time, mass, and charge follow at the projection boundary. They coincide with what physics calls the Planck scale, but they are defined from the substrate, not from SI constants. The SI values are verification, not foundation.

One open question stays explicit: why does the projection carry what the observer calls units of action --- energy times time --- and not some other combination? Answering it means deriving the observer's dimensional categories from the substrate. That derivation does not yet exist. It is a primary target, carried in Unresolved Questions.

\emph{Status: Established in v0.5.}

\subsection{Emergent Properties}
\label{sec:org384c46f}

The Medium distinguishes three categories of emergent property, organized by the level at which they arise.

\textbf{Substrate emergent.} These properties arise directly from \(\iota\)-combinations, prior to any observer: mass (corresponding to \(N\) as a first approximation --- the count of \(\iota\) events in a configuration), charge (corresponding to the net distinction character of the \(\iota\) configuration), spin (arising from relational closure of the configuration's encounter pattern within the web), the quantum vacuum (sub-threshold substrate combinations, appearing at the view layer as vacuum fluctuations --- measurable consequences such as the Casimir effect and Lamb shift are open derivations), the Hamiltonian (projected from the substrate combination structure), spacetime (as the projection framework constructed by embedded observers), and dimensionality (as the coordinate structure observers impose on substrate relationships).

\textbf{View layer emergent.} These properties arise from interactions between stable configurations within the spacetime projection: the fundamental forces, quantum fields, and the full complexity of observable physical reality. This is the domain of existing physics. The Medium makes no claims in this domain beyond identifying it as downstream.

\textbf{Unassigned.} The following have not yet been placed with confidence: neutrino oscillation and dark matter. Dark energy dissolves as a projection artifact of misread redshift. Color charge dissolves as a labeling convention with no substrate reality. Gravity has a candidate direction: mass is a stable stochastic imperfection in the \(\iota\) packing; gravity is the strain field propagating through the medium toward ground-state restoration. The current limit of assignment reflects the limit of present understanding, not the limit of The Medium.

\subsection{Mass, Charge, and Spin as Projection Signatures}
\label{sec:orgdab1f59}

Standard physics treats mass, charge, and spin as independent intrinsic properties of particles. The Medium identifies them as three projection signatures of a single underlying combinatorial structure.

Mass corresponds to \(N\) as a first approximation: the total count of \(\iota\) events in a stable configuration. A heavier particle is a configuration with a larger \(N\). The specific values --- why the muon is 207 times the electron's mass, why three generations exist at all --- are not yet derived from substrate structure.

Charge corresponds to the net distinction character of the \(\iota\) configuration --- the combination of its \(\iota\) events' relational modes. A configuration in which one distinction character dominates projects to positive charge; the opposing character projects to negative charge; exact balance projects to charge neutrality. Charge is not a property added to a particle. It is the distinction character of its substrate configuration.

Spin describes how a stable configuration's encounter pattern closes within the relational web. It is not rotation. The particle does not spin. The word is retained from the observation, but the property is relational, not mechanical.

A configuration whose encounter pattern returns to its original relational state after one traversal of the web projects as integer spin. A configuration requiring two traversals projects as half-integer spin. The distinction is not internal to the configuration alone. It is in how the configuration's encounter pattern sits within the surrounding web.

The specific derivation of which encounter patterns produce which spin values (0, 1/2, 1, 3/2, 2) remains an open problem. The claim that spin is a property of relational closure is current.

\subsection{Observer and Observed}
\label{sec:orga280e4d}


Standard quantum mechanics imports a Newtonian posture. A detached external entity measuring a passive system. That posture breaks down at the quantum scale. It was never properly grounded anywhere.

In The Medium there is no outside. Every interaction is between \(\iota\)-level processes within the Medium.

The observed is any process whose boundary state is altered by interaction. The observer is any process that retains a structural consequence of that interaction — registration. The asymmetry dissolves. Registration is mutual and structural. Not performed from a privileged external position.

A note on vocabulary. The word observation carries two incompatible meanings that the framework must keep distinct. In QM, observation means the interaction that produces a definite outcome — no consciousness required, purely structural. In ordinary language, observation means the experience of a perceiving entity. The Medium uses registration for the first and observer for the second. Registration is the structural event at the projection boundary. The observer is the Level 3 embedded configuration constructing a dimensional account of what the relational web presents. These are different things at different levels. Conflating them is the source of a century of confusion about the measurement problem.

Superposition is not a rapid toggling between states — interference patterns from the double slit rule that out definitively. It is the formalism's acknowledgment that before interaction, definite states do not exist. The measurement problem has gone unsolved for a century because the observer/observed distinction was never grounded. The Medium dissolves the distinction at the substrate level. What physics calls measurement is registration. What physics calls observer is any process capable of retaining a structural consequence of boundary contact.

\subsection{Co-Manifestation and Entanglement}
\label{sec:orgcd23f22}

At the substrate level, two \(\iota\) events producing the same \(\phi\) character share a state with no spatial relationship between them. There is no distance at the substrate level. There is no \emph{between}.

When an observer embedded in differentiated structure encounters two events that are correlated regardless of their spatial separation, it calls this entanglement. Entanglement is the observer's projection of co-manifestation \(x \sim y\) into a framework that expects spatial locality to govern correlation. The apparent mystery of entanglement --- correlation without causal connection across space --- dissolves at the substrate level, where space does not yet exist and co-manifestation requires no connection at all.

\subsection{The Stochastic Character --- Observer's Description}
\label{sec:org4de6abd}

At the substrate level, \(\iota\) events are not driven. No prior state determines which relational node the Medium next distinguishes within itself. Each act of self-distinction produces a character --- +1 or -1 --- with no preference and no guarantee. An observer describing this stochastic character mathematically recognizes it as a Bernoulli process --- a binary outcome with associated probability \(p(\iota)\) for each \(\iota\) event. This is the observer's mathematical framework imposed on the substrate's stochastic character. The substrate does not know about Bernoulli. Distinction simply occurs or does not.

The full statistical mechanics of the ground state --- entropy, temperature, distributions --- emerges on the observer side as the projection of the substrate's undifferentiated stochastic character into the observer's thermodynamic framework.

\subsection{The General Principle}
\label{sec:orga16a0e8}

The presence of dimensional constants in the equations of physics --- \(h\), \(c\), \(G\), \(k\) --- is the signature of projection. Each constant marks a place where the observer's dimensional framework has been imposed on a substrate relationship that has none. In the substrate's own terms, those constants do not appear.

The program implied by the symbol table: for each row relating a substrate quantity to an observer quantity, derive the dimensional form of the observer quantity from the substrate quantity and a formal account of the projection boundary. If that program succeeds, every fundamental constant of physics will have been given a substrate derivation --- and the mystery of their values resolved, not by finding deeper constants, but by understanding that they were never fundamental to begin with.

\subsection{On Existing Physics}
\label{sec:orgd66d526}

The existing mathematical frameworks of physics are valid descriptions of differentiated structure behavior within their respective domains. They are not wrong. They are downstream. They are precise and consistent descriptions of the observer's projection of the substrate.

The Medium framework does not replace them. It identifies the substrate from which they project --- and defines the program by which their dimensional constants can, in principle, be derived rather than assumed.

Three features of existing physics that appear to be gaps in The Medium are not gaps. They are artifacts of the mathematical starting point those frameworks adopted.

\textbf{The complex structure of quantum mechanics.} Quantum mechanics requires complex-valued wave functions --- amplitudes are complex numbers, the Schrödinger equation contains an explicit imaginary unit. This was not derived from first principles. It was adopted: Schrödinger borrowed wave mathematics from classical physics; the complex Hilbert space formulation was chosen because it worked. The Medium does not need to derive complex numbers. It needs to derive the phenomena --- interference, superposition, the uncertainty relations --- from substrate structure. If those phenomena emerge from co-manifestation and relational closure, the observer will find that complex Hilbert space is a useful mathematical description of what they are seeing. That is the description, not the ground. The complex structure is a view-layer formalism built on a wrong starting point. The phenomena are real. The mathematics is downstream.

\textbf{Renormalization.} Quantum field theory on continuous spacetime produces infinities in its calculations. These are absorbed into redefined parameters through renormalization --- a mathematically consistent but physically unsatisfying procedure. The infinities arise because the theory assumes no minimum length scale: arbitrarily short wavelengths contribute to loop integrals. The Medium has a minimum scale: one \(\iota\) event is the irreducible unit of distinction. Below that there is nothing to integrate over. The UV divergences that require renormalization do not arise. The problem dissolves because the substrate removes the assumption that generated it.

\textbf{The Schrödinger equation.} If the Hamiltonian is derived from the substrate combination structure --- \(\mathrm{proj}(\Kappa)\) in the scalar account, pending mode-structure reformulation --- and time is view-layer emergent, then the Schrödinger equation \(i\hbar \partial\psi/\partial t = \hat{H}\psi\) is a view-layer description --- the observer's account of how configurations evolve in projected time under the derived Hamiltonian. It is not a substrate input. It is downstream of the substrate structure from which the Hamiltonian emerges.

\textbf{The uncertainty principle.} Heisenberg's uncertainty principle --- \(\Delta x \Delta p \geq \hbar/2\) --- is often presented as a fundamental postulate. It is not. It is a mathematical consequence of Fourier transforms: a wave with perfectly definite frequency extends infinitely in space; localizing a wave requires superposing many frequencies; precise position means indefinite momentum and vice versa. That is a mathematical relationship, not a physical constraint imposed from outside. In The Medium, position and momentum are both view-layer constructs. Neither exists at the substrate level. When the observer projects substrate structure into a coordinate system and assigns position and momentum, the Fourier relationship between the two conjugate constructs is inherited from the wave structure of the projection itself. The uncertainty principle follows automatically. It is not a substrate feature requiring separate derivation. The stochastic character of \(\iota\) events --- no prior state determines the next --- is the substrate ground of the uncertainty. The specific form \(\Delta x \Delta p \geq \hbar/2\) is the view-layer expression of that ground, shaped by the Fourier relationship between the position and momentum bases the observer constructs.

\textbf{The information paradox.} Hawking showed that black holes emit thermal radiation carrying no information about infalling matter. If a black hole evaporates completely, information appears to be destroyed. This contradicts quantum mechanical unitarity --- unitary time evolution requires information to be conserved. The paradox has been central to theoretical physics for fifty years. The Medium reframes it. Unitarity is a property of the Schrödinger equation. The Schrödinger equation is view-layer --- the observer's account of how configurations evolve in projected time. It is not a substrate input. The substrate is stochastic: no prior state determines the next \(\iota\) event. There is no substrate-level requirement for deterministic reversibility. Unitarity is a view-layer property that holds for isolated systems that remain within the view layer. When a black hole forms and evaporates, the configurations that produced it dissolve back into unconstrained \(\iota\) events below the projection threshold. From the view layer this looks like information loss. From the substrate it is the stochastic ground returning to its undifferentiated state. Information is not destroyed --- it retreats below view-layer resolution. A process that crosses below the projection boundary was never subject to view-layer unitarity to begin with. The paradox is reframed, not dissolved. Whether that reframing is sufficient to close it is an open question.

\textbf{Bell's theorem.} Bell's theorem proves that no local hidden variable theory can reproduce quantum mechanical predictions. Experiments confirm quantum mechanics. Local hidden variable theories are ruled out. The Medium's substrate may appear to be a hidden variable theory --- \(\iota\) events are below the observational floor, not directly measured, determining the characters that project to observable physics. That pattern is what hidden variable means. The Medium is not a local hidden variable theory. Bell's theorem requires locality: influences cannot propagate faster than \(c\) between spacelike-separated events. The Medium's substrate has no \(c\) --- \(c\) is view-layer emergent. The substrate has no spatial separation --- space is view-layer. Co-manifestation is explicitly non-local: two \(\iota\) events sharing a state with no spatial relationship between them. There is no distance at the substrate level. Bell's locality assumption does not hold at the substrate level. The theorem's premises do not apply. The Medium is a non-local substrate theory. The locality that Bell requires is a view-layer artifact that does not exist at the ground floor. The apparent mystery of entanglement --- correlation without causal connection across space --- is not mysterious at the substrate level. Space does not yet exist there. Co-manifestation requires no connection across it.


\section{PART IV: COSMOLOGICAL IMPLICATIONS}
\label{sec:org67d0784}

The following develops four interconnected consequences that follow directly from the foundational structure of The Medium. Each addresses a major puzzle in standard cosmology. In each case the puzzle dissolves rather than being solved --- it was generated by a category error that The Medium does not commit: treating view-layer artifacts as substrate-level facts.

\subsection{Entropy as Emergent Accounting}
\label{sec:org7996cc6}

Standard cosmology inherits a puzzle from Boltzmann: if the second law of thermodynamics is statistical rather than fundamental, why was entropy low at the beginning? The standard response bundles this into initial conditions and calls it the past hypothesis. The Big Bang is asserted to have been an extraordinarily low-entropy state. This assertion is not derived. It is required by the model and inserted by hand.

The Medium does not accept this framing, and not merely because it rejects the Big Bang as an explanation. The more fundamental problem is that the framing assigns entropy as a global intrinsic property of the universe --- a quantity that belongs to the cosmos as a whole and increases over cosmic time. This is a category error.

Entropy, as Boltzmann and later Shannon showed, is a measure of how much an observer cannot track about a system. It is not a property of the system alone. It lives in the gap between an observer with limited reach and a system doing more than the observer can follow. It always depends on where the observer draws the line --- which differences count and which do not.

The universe as a whole has no external observer. Assigning it a global entropy value requires an observer standing outside it with a complete accounting of all microstates, which is incoherent. What cosmologists call the entropy of the universe is always the entropy of the observable universe as seen from a particular embedded vantage point with particular resolution limits. It is local. It is relational. It is view-layer all the way down.

In The Medium, the substrate --- the \(\iota\)-combinatorial structure --- has no thermodynamic direction. There is no privileged low-entropy initial condition because there is no initial condition at the substrate level in any cosmologically meaningful sense. The apparent entropy gradient that standard cosmology treats as a cosmic fact is a view-layer phenomenon: an artifact of embedded observers constructing temporal sequences and tracking accessible macrostates from within a particular slice of configuration space.

Entropy is therefore classified in The Medium as view-layer emergent. It has no substrate analog.

\subsection{Time as Projection and Its Consequences for Frequency}
\label{sec:org7cefc4b}

The Medium establishes that time is not a substrate-level property. Temporal directionality --- the sense that events occur in sequence, that causes precede effects, that entropy increases in one direction --- emerges at the view layer as observers at a particular scale construct ordered accounts of substrate transitions.

Time is real for the embedded observer. Thermodynamically, experientially, irreversibly real. Dismissing it as illusion is philosophical evasion. But real-for-the-observer and necessary-for-the-universe are different claims entirely. The universe's processes do not require time to coordinate. Gravity does not consult a clock. The gravitational relationship between earth and moon is not a sequence of instructions passed back and forth --- it is a continuous configuration of the substrate, of which the 1.3-second propagation delay at the view layer is the observer's local account. Physics wrote the observer's experience of sequence into the equations as a global feature of reality. The equations keep resisting it. General relativity is time-symmetric. Physics's deepest equations keep coming out with no time in them. Time is emergent downstream of process. Not upstream of it.

This has a direct consequence for the concept of frequency. Frequency is defined as cycles per unit time. If time is view-layer emergent, then frequency is also view-layer emergent. It is not a substrate-level property of a propagating pattern. It is a measurement that an observer constructs by counting transition events against a locally projected time reference.

Two observers projecting different temporal frames --- because they sit at different configuration depths within the substrate, or because their local \(\iota\)-configuration densities differ --- will assign different frequency values to the same arriving pattern. This is not a Doppler effect. It is not recession. It is a structural consequence of time being a local projection rather than a universal background.

\subsection{Redshift as Multi-Layer Projection Artifact}
\label{sec:orgab3c5b2}

Standard cosmology interprets cosmological redshift as evidence of recession velocity: distant galaxies are moving away, space between them and us is expanding, and their light is stretched to longer wavelengths in proportion to the expansion. This interpretation underlies the Big Bang model, Hubble's law, and the entire framework of an expanding universe with a singular low-entropy origin.

The Medium offers a more parsimonious account that requires no expansion, no recession, and no singular origin --- while accepting the observational data entirely.

Within The Medium, what we call a photon is a stable propagating \(\iota\)-configuration pattern --- a structure that persists and moves through successive co-manifestation relationships across configuration space. What we call speed is a view-layer quantity: a ratio of projected distance to projected time constructed by the observer. It is not a geometric velocity through space. It is not a substrate-level rate. The substrate has no rates --- no shared temporal framework exists against which transitions could be counted.

Redshift arises from three stacked projection effects.

\emph{First: the emitting configuration.} A source at some region of the substrate produces a transition pattern at a rate determined by local \(\iota\)-configuration density. That rate is what the source's local view layer calls frequency.

\emph{Second: traversal through configuration depth.} The pattern propagates across substrate regions that are not uniform. Configuration density varies across the substrate --- not directionally, not as expansion, but structurally, as a feature of the combinatorial landscape. As the pattern traverses this varying density, the relationship between transition events and projected distance is not constant.

\emph{Third: the receiving configuration.} The observer receives the arriving pattern and measures its frequency against a locally projected time reference. That time reference is constructed from the observer's local \(\iota\)-configuration density, which differs from the emitter's. The frequency the observer assigns is a ratio constructed entirely within the view layer, against a local clock that has no substrate-level universality.

The redshift that results from this three-layer process is not evidence that the source is receding. It is evidence that the emitting and receiving configurations sit at different depths within a non-uniform substrate, and that the projection relationship between them is asymmetric. What standard cosmology interprets as Hubble's law --- a proportional relationship between distance and recession velocity --- is, in The Medium, a projection artifact of sampling along a particular observational cone from a single embedded vantage point through a substrate with a particular large-scale configuration structure.

This account creates an obligation. If redshift is configuration-depth asymmetry rather than recession, the linearity of Hubble's law requires explanation. A clean proportional relationship between redshift and distance is a very specific result. It implies that the large-scale \(\iota\) configuration structure has a particular gradient --- regular enough to produce linearity along any observational cone from any embedded vantage point. Why the substrate's large-scale structure produces this regularity is an open derivation. It is named here as a primary target, not dismissed. A projection artifact that produced a random or irregular redshift-distance relationship would be immediately falsified. The linearity is a constraint the iota account must satisfy.

The mature and fully structured galaxies observed at high redshift --- including findings from JWST that are unexpected under standard galaxy formation timelines --- are consistent with this account. They do not violate a law of nature. They strain the galaxy formation models layered on top of the Big Bang framework, not the framework's core timeline. In The Medium, structural complexity at any configuration depth is not constrained by elapsed cosmic time, because cosmic time is itself a projection.

Dark energy dissolves in this account. It was not derived. It was a fix. If redshift is configuration-depth asymmetry rather than recession, there is no expansion to accelerate. And if nothing is accelerating, there is nothing for dark energy to explain. It is gone. One honest question is left: why don't the small substrate fluctuations add up to a pull that gravity would notice? I don't have that answer yet.

\subsection{The Big Bang as Projection Error}
\label{sec:orgf04fb5f}

The Big Bang model rests on three foundational assumptions, each of which The Medium addresses directly.

First: that cosmological redshift indicates recession velocity, implying an expanding universe with a singular dense origin. Second: that entropy has a global cosmic direction, requiring a special low-entropy initial state. Third: that time provides a universal background against which an origin event --- a moment of creation --- is a coherent concept.

The Medium has shown that all three assumptions are view-layer artifacts. Redshift is configuration-depth asymmetry misread as velocity. Entropy is a local observer relationship with no substrate reality. Time is a local projection with no substrate universality.

Remove those three assumptions and the interpretation built on them loses its foundation. The observational data remain --- including the cosmic microwave background and the abundance ratios of hydrogen, helium, deuterium, and lithium predicted by Big Bang nucleosynthesis. The Medium does not yet account for those. They are carried as open problems in the Unresolved Questions. What does not survive is the interpretive framework erected on top of the data --- the claim that redshift requires recession, that entropy requires a special initial state, that time requires a singular origin.

This is not a disproof. A disproof accepts a model's foundations and finds an internal flaw. The Medium does something else. It questions the foundations themselves --- recession, a global entropy direction, a universal clock. The data are not in dispute. The interpretation built on them is.

The universe, in The Medium, has no beginning in any substrate-meaningful sense. It has no privileged center, no singular origin, no global clock, and no cosmic entropy gradient. What it has is a substrate of extraordinary combinatorial richness from which structure, observers, and the local projections we call time, space, and thermodynamics all emerge --- consistently, without residue, and without requiring any special initial conditions.

\subsection{The Flatness Problem as Projection Artifact}
\label{sec:org87cac8d}

Standard cosmology requires the universe to be extraordinarily flat --- far flatter than it has any right to be. Flatness is unstable under standard expansion: small deviations grow over time. For the universe to be this flat today, initial conditions had to be extraordinarily fine-tuned. Inflation was introduced partly to solve this --- exponential expansion drives any initial curvature to negligible levels.

The Medium dissolves the problem without inflation.

Space is view-layer emergent. The geometry of space is constructed by embedded observers from the relational web of co-manifestation. There is no pre-existing spatial geometry to be curved or flat. The relational web has a large-scale structure determined by the statistical properties of co-manifestation across the substrate. The substrate has no preferred direction --- direction is view-layer. The projected geometry therefore inherits no preferred curvature. Flatness is not a fine-tuned initial condition. It is the generic result of projecting a substrate with no preferred spatial orientation. Any large-scale curvature would require a preferred direction in the substrate. No such direction exists.

The flatness problem is generated by assuming space is a pre-existing container that could have started curved. Remove that assumption and the fine-tuning dissolves. The universe is flat because the substrate from which space is projected has no mechanism for producing preferred curvature.


\subsection{The Horizon Problem as Projection Artifact}
\label{sec:org9dffdea}

Standard cosmology's horizon problem starts with the cosmic microwave background --- the faint afterglow left over from the early universe, spread across the whole sky. Measured in every direction, its temperature comes back the same. That sameness is an observation, not an assumption. The puzzle: two regions of the sky far enough apart could never have traded a light signal in the time available. They never touched. Yet their temperatures match. Inflation patches this by stretching one small connected region out to cover the whole sky.

The Medium dissolves the uniformity without inflation.

Causal horizons require two things: a speed limit (\(c\)) and a time since origin. Both are view-layer in The Medium. \(c\) is a ratio of projected distance to projected time. The Big Bang timeline is a projection artifact. No substrate-level causal horizon exists because there is no substrate-level \(c\) or elapsed time.

The uniformity of the CMB is the projection of the substrate's statistical homogeneity. The substrate has no preferred location. \(\iota\) events are statistically uniform across the relational web. That statistical uniformity projects as uniform temperature. No causal contact was required because causality is view-layer. The two apparently disconnected regions were never disconnected at the substrate level --- spatial separation does not exist there.

The uniformity is handled. The tiny temperature ripples that later grew into galaxies are not. They need a substrate account of which configurations make which ripples, and I don't have one. It is carried in the CMB entry of Unresolved Questions.

\section{PART V: WHAT THE MEDIUM NEEDS TO PROVE}
\label{sec:orgc7a36d9}


Not a physical theory with testable predictions in the conventional sense. Not a mathematical framework with full derivational power yet.

But not untestable.

The framework makes one class of indirect prediction that follows directly from its central claim. If particle properties are relational outcomes of substrate combinatorics rather than intrinsic properties of the particles themselves, then those properties should show systematic context-dependence that the intrinsic model cannot explain. The measurement context engages the configuration's relational modes. Different contexts engage different modes. Properties that are intrinsic should be invariant across measurement contexts. Properties that are relational should not be.

This is a real, testable difference. If a property is relational, it should shift with the measurement context --- and the shift should be larger for heavier, more complex particles, smaller for simple ones. Fixed-property physics predicts no such shift. Measure it and the two views come apart. If the shift is not there, my account is wrong.

The framework does not predict new particles. It predicts structure where standard physics expects none.

Three things the framework needs to prove.

First: that the primitives are genuinely primitive. That occurrence, registration, boundary do not smuggle in hidden assumptions. That they do not secretly depend on Newtonian concepts.

Second: that existing frameworks require something like these primitives to be coherent. That the measurement problem, the time problem, the observer confusion all point to a gap sitting exactly where The Medium works. Three frameworks in particular are candidates for derivation from the substrate.

Planck bits are what encounter characters look like to an observer who cannot see below the registration level. Digital physics assumes the bit as primitive. The Medium shows where it comes from. The path from encounter to registered character to bit is the derivation. It is not complete. The direction is clear.

Loop quantum gravity makes spacetime emergent from spin networks --- nodes and edges that are atoms of geometry. The Medium's encounter web is already that relational structure. Stable encounter configurations at stability extrema are the nodes. The encounter relationships between them are the edges. LQG's spin foam --- the four-dimensional history of spin networks --- is the observer's projection of encounter patterns across what it calls time. The derivation would show that spin networks are what the encounter web looks like when an observer imposes a spatial coordinate system on it.

String theory places vibrating strings at the foundation. Different vibration modes correspond to different particles. In The Medium, different encounter configurations at stability extrema correspond to different particles. The correspondence is directionally plausible. String theory presupposes spacetime and requires ten or eleven dimensions. The derivation would need to show those dimensions emerging from encounter structure. That is the longest path of the three.

None of these derivations is complete. All three are named targets. Showing that existing frameworks are derivable from the substrate is the substance of this second criterion.

Third: that the framework is internally consistent and generative. That thinking from these primitives produces coherent consequences. That it opens questions existing frameworks could not formulate.

The model is the 1905 paper on special relativity. Not complex. Precise. Spare. The argument careful. The boundaries clearly stated.

That is achievable. And it is enough.

\subsection{Prior Frameworks: What They Reached and Where They Stopped}
\label{sec:orgf9b3368}

Three frameworks reach toward the same ground and stop short. A fourth is instructive in a different way.

\textbf{Planck bits.} Wheeler's \emph{It from Bit}, developed further by Fredkin and Wolfram, places a binary digit at the foundation. The bit carries its state intrinsically. The Medium shows where the bit comes from: it is the registered distinction character of one \(\iota\) event. Digital physics correctly identifies the information structure of reality. It misidentifies that structure as the ground floor.

\textbf{Loop quantum gravity.} Smolin, Rovelli, and collaborators developed a background-independent approach to spacetime in which geometry emerges from a relational network --- spin networks of nodes and edges, with no pre-existing spatial container. This is a genuine step toward the same territory. The relational web of The Medium is structurally similar: position defined by co-manifestation, adjacency defined by directness of relation, spacetime emergent from the network rather than assumed. The difference is one level. LQG assumes the relational network as its starting point. The Medium derives the relational network from something more primitive: the \(\iota\) and its intrinsic attraction. That is the step LQG does not take.

\textbf{String theory.} Different vibrational modes of a fundamental string correspond to different particles. The Medium's different \(\iota\) configurations at stability extrema correspond to different particles. The structural parallel is real. String theory presupposes spacetime and requires ten or eleven dimensions. The derivation from The Medium would need to show those dimensions emerging from the relational mode structure. That is the longest path of the three.

\textbf{The computational model.} A recurring analogy in foundational physics holds that reality is computational --- that the universe processes information according to rules, and physical phenomena are the outputs. Smolin, Wheeler, and Wolfram each approach this framing from different directions. The analogy is instructive but fails at a critical point. Computation requires rules imposed from outside: a program, a set of instructions, something that specifies what the computation does. Even cellular automata have rules baked in before the first step. The relational web of The Medium has no imposed rules. The structure is what intrinsic \(\iota\) attraction does at scale. Nobody programmed the packing geometry. It falls out of what the \(\iota\) is. This is not computation. It is emergence. The distinction matters: a computational universe requires an external rule-giver, which simply relocates the foundational question. An emergent universe does not.

\section{PART VI: DERIVATIONS}
\label{sec:org3ab4e00}


\subsection{Derivation I: Planck Bits from the Substrate}
\label{sec:org06cb493}

Digital physics --- Wheeler's \emph{It from Bit}, Fredkin, Wolfram --- claims the bit as primitive. A binary value, 0 or 1, is the most fundamental thing. Everything physical arises from it. No explanation is offered for where the bit comes from. It simply is.

The Medium derives it.

\textbf{Step 1.} \(\iota\) occurs. The Medium distinguishes within itself. That is the event. It is not driven. No prior state determines it. It simply occurs.

\textbf{Step 2.} The distinction has two aspects: +1 or -1. Not chosen from outside. Not determined in advance. The character of the distinction itself. Without difference there is no distinction. The character is what distinguishing is.

\textbf{Step 3.} The character is registered in \(\phi\) for that \(\iota\) event. The field records it. One distinction. One binary value.

\textbf{Step 4.} At the projection boundary, one \(\iota\) event corresponds to one Planck-scale event. This is the threshold at which the substrate's dimensionless distinction acquires the dimensional form the observer calls action. One \(\iota\) event. One Planck unit. One registration.

\textbf{Step 5.} What the observer records at the Planck scale is one binary value: +1 or -1.

That binary value is the bit.

A Planck bit is not a primitive. It is the registered distinction character of one \(\iota\) event, as seen from the observer's position. The bit presupposes the act of distinction. Digital physics correctly identifies the information structure of reality. It misidentifies that structure as the ground floor.

\textbf{Consequences that follow without additional derivation:}

The stochastic character of bits --- no prior state determines the next --- follows directly from \(\iota\) events not being driven. No prior state determines where or when the Medium next distinguishes within itself.

The minimum size of a Planck bit --- information cannot be compressed below one Planck unit --- follows from \(\iota\) being the unit of distinction. Below one distinction event, there is nothing to register. The Planck scale is not arbitrary. It is the projection boundary of the substrate's irreducible unit.

The adjacency constraint explains why directly co-manifested identical bits produce no new information. Two \(\iota\) events of identical distinction character in direct co-manifestation produce no new distinction. This is not a rule imposed on the bits. It is a logical consequence of what distinction is.

\textbf{What is not yet derived:}

Why one \(\iota\) event projects to exactly \(\hbar\) of action --- energy multiplied by time --- and not some other dimensional combination. This remains in unresolved questions. The derivation above stands without it.

\subsection{Derivation II: Standard Particles as Stable \(\iota\) Configurations}
\label{sec:orgccd2e84}

\emph{Note: This derivation uses the adjacency constraint as working version and names stability extrema as the target condition. The formal combination rules governing \(\iota\) relational modes are an open derivation. The particle table derivation is provisional until those rules are established. The direction is current; the formal path is open.}

The Planck bit derivation has an immediate consequence. If a Planck bit is the observer's view-layer name for a registered \(\iota\) distinction character, then what an observer calls a standard particle is a stable \(\iota\) configuration --- a pattern of distinction characters that has achieved a stability extremum in the relational web.

The properties of that configuration are what the pattern projects to the observer.

Mass corresponds to \(N\) as a first approximation: the count of \(\iota\) events in the configuration. A heavier particle is a larger pattern. The lightest particles are the threshold cases --- the minimum number of \(\iota\) events that can achieve stability under the combination rules. The specific mass values and the three-generation hierarchy are not yet derived.

Charge corresponds to net distinction character. A pattern in which one character dominates projects as charged. Exact balance projects as neutral.

Spin corresponds to relational closure. The pattern either closes in one traversal of the web or two. One traversal: boson. Two traversals: fermion.

\textbf{The combination rules shape the particle table.}

Identical adjacent distinction characters produce nothing. Only patterns that vary can build. This sharply limits which configurations are possible. The number of stable particles is not arbitrary. It is constrained by which patterns satisfy the combination rules. The particle table may be exactly the set of patterns the rules permit. The adjacency constraint is the working version of those rules.

\textbf{The threshold cases:}

The electron is the minimum stable configuration with net character imbalance and half-integer closure. Possibly the simplest stable fermion the combination rules allow.

The photon does not sit at a stability extremum. It is a propagating alternation of distinction characters --- the combination rules satisfied but no bound state achieved. The observer assigns it speed \(c\) --- a view-layer ratio, not a substrate property. Why projection fixes that ratio at \(c\) and not some other value is an open derivation.

The neutrino approaches the stability threshold from below. Near-zero mass. Balanced character. The minimum fermion-closing configuration.

Quarks cannot achieve stability alone under the combination rules. Stability requires combination. The structural impossibility of isolated quark stability is a substrate consequence of the combination rules. The energy-dependent behavior of the strong force --- asymptotic freedom, running coupling --- is a view-layer QCD result. These are distinct claims at different levels.

\textbf{The open derivation:}

Which specific \(\iota\) patterns --- satisfying the combination rules, achieving a stability extremum --- correspond to each observed particle? That is the derivation. It begins with the lightest threshold cases and works upward through the particle table. The framework predicts that the table is finite and complete. The combination rules permit exactly the particles that exist. No derivation yet confirms this. It is the primary structural task of the next stage.

\subsection{Derivation III: \(G\) from the Vacuum Floor and Hubble Rate}
\label{sec:org6523417}

\emph{Note: This derivation identifies the structural form of \(G\) in Medium Theory terms and produces 94--97\% numerical agreement with the measured value. The remaining shortfall is explained and carries a testable equivalent prediction. A fully \(G\)-independent derivation requires the formal strain field equations, which are open work.}

\subsubsection{The Setup}
\label{sec:orgd45787d}

Gravity in The Medium is the strain field through the \(\iota\) plenum as it corrects toward ground state. At the view layer, this strain field projects as the Newtonian gravitational potential \(\Phi\) satisfying Poisson's equation:

\begin{equation}
\nabla^2 \Phi = 4\pi G \rho
\end{equation}

where \(\rho\) is mass density and \(G\) is the coupling constant --- the medium's compliance to gravitational perturbations. \(G\) is a projection seam constant: the observer's SI measurement of a substrate coupling ratio that has no dimensions at the substrate level.

The question is which substrate quantities set \(G\).

\subsubsection{The Two Scale-Setting Quantities}
\label{sec:org19b99a1}

The medium at cosmological scale has two measurable characteristic quantities:

\(\rho_\Lambda\) --- \emph{the vacuum floor.} The minimum projected \(\iota\) activity: the medium's rest-state energy density. This is the dark energy density, the term the observer associates with the cosmological constant \(\Lambda\). Below \(\rho_\Lambda\), no substrate activity projects to the view layer. It sets the stiffness of the medium --- the resistance to being perturbed away from ground state.

\(H_0\) --- \emph{the Hubble rate.} The large-scale \(\iota\) configuration-depth gradient at the current epoch, which the observer reads as Hubble's law. In The Medium, \(H_0\) is not a recession rate but the characteristic frequency at which large-scale \(\iota\) density gradients manifest at the projection boundary.

Both are inputs to the strain field coupling. \(G\) must be constructible from them.

\subsubsection{The Dimensional Argument}
\label{sec:orgd121d59}

\(G\) has SI units \(\mathrm{m^3\,kg^{-1}\,s^{-2}}\).

The available quantities and their units:
\begin{align*}
[\rho_\Lambda] &= \mathrm{kg\,m^{-3}}\\
[H_0] &= \mathrm{s^{-1}}\\
\left[\frac{H_0^2}{\rho_\Lambda}\right] &= \frac{\mathrm{s^{-2}}}{\mathrm{kg\,m^{-3}}} = \mathrm{m^3\,kg^{-1}\,s^{-2}}
\end{align*}

\(H_0^2/\rho_\Lambda\) is the unique combination of the vacuum floor and Hubble rate that carries the units of \(G\). No other combination of these two quantities works.

The factor \(4\pi\) in Poisson's equation is geometric --- it arises from Gauss's theorem in three spatial dimensions (the surface integral over a sphere in 3D is \(4\pi r^2\)). It is not a free parameter; it is fixed by the dimensionality of the projection.

\subsubsection{The Result}
\label{sec:orge6b9258}

\begin{equation}
\boxed{G = \frac{H_0^2}{4\pi\rho_\Lambda}}
\end{equation}

\emph{Physical interpretation.} \(\rho_\Lambda\) is the stiffness: a higher vacuum floor means a stiffer medium, less compliant to gravitational perturbations, and therefore a smaller \(G\). \(H_0^2\) is the dynamic factor: a higher large-scale oscillation rate produces stronger coupling. The ratio \(H_0^2/\rho_\Lambda\) is the medium's compliance --- how readily the strain field responds to a mass imperfection. \(G\) is that compliance scaled by the 3D solid-angle factor \(1/(4\pi)\).

Matter is excluded from \(\rho_\Lambda\) because matter is the imperfection --- the perturbation the strain field responds to. Matter contributes to the total stress-energy budget but not to the ground-state stiffness of the medium. Only the vacuum floor sets the compliance.

\subsubsection{Numerical Verification}
\label{sec:org2ff3c19}

Using Planck 2018 cosmological parameters (\(H_0 = 67.36\) km/s/Mpc, \(\Omega_\Lambda = 0.685\)):

\begin{align*}
H_0 &= 2.183 \times 10^{-18}\ \mathrm{s^{-1}}\\
\rho_\Lambda &= \Omega_\Lambda \times \rho_c = 5.836 \times 10^{-27}\ \mathrm{kg\,m^{-3}}\\[6pt]
G_{\mathrm{derived}} &= \frac{(2.183\times10^{-18})^2}{4\pi\,(5.836\times10^{-27})}
= 6.498 \times 10^{-11}\ \mathrm{m^3\,kg^{-1}\,s^{-2}}\\[4pt]
G_{\mathrm{measured}} &= 6.674 \times 10^{-11}\ \mathrm{m^3\,kg^{-1}\,s^{-2}}\\[4pt]
\text{Agreement} &= 97.4\%
\end{align*}

With WMAP9-era parameters (\(H_0 = 68.0\), \(\Omega_\Lambda = 0.71\)): agreement = 94.0$\backslash$%.

\subsubsection{The Analytic Result}
\label{sec:org526d815}

The formula has an exact closed form in terms of \(\Omega_\Lambda\). Substituting \(\rho_\Lambda = \Omega_\Lambda\,\rho_c\) and the Friedmann equation \(\rho_c = 3H_0^2/(8\pi G)\):

\begin{align*}
G_{\mathrm{derived}} &= \frac{H_0^2}{4\pi\,\Omega_\Lambda\,\rho_c}
= \frac{H_0^2}{4\pi\,\Omega_\Lambda} \cdot \frac{8\pi G}{3H_0^2}
= \frac{2G}{3\Omega_\Lambda}
\end{align*}

Therefore:

\begin{equation}
\frac{G_{\mathrm{derived}}}{G} = \frac{2}{3\Omega_\Lambda}
\end{equation}

This is an exact result. It means the formula \(G = H_0^2/(4\pi\rho_\Lambda)\) is precisely equivalent to the prediction:

\begin{equation}
\Omega_\Lambda = \frac{2}{3} \approx 0.667
\end{equation}

Observed: \(\Omega_\Lambda \approx 0.685\) (Planck 2018). Discrepancy: 2.7$\backslash$%.

\subsubsection{Interpreting the Shortfall}
\label{sec:orga0c32f1}

The 2.7$\backslash$% discrepancy is the matter contribution. The Friedmann equation governs the full stress-energy budget including matter and radiation:

\[
H_0^2 = \frac{8\pi G}{3}(\rho_\Lambda + \rho_m + \rho_r)
\]

The Medium formula uses only \(\rho_\Lambda\) in the denominator. The matter density \(\rho_m \approx 0.315\,\rho_c\) contributes to the Friedmann sum but not to the vacuum floor stiffness. The current measurement of \(\rho_\Lambda\) in SI units uses \(G\) (through \(\rho_c = 3H_0^2/8\pi G\)), which slightly inflates the apparent \(\rho_\Lambda\) and suppresses \(G_{\mathrm{derived}}\).

When the vacuum floor is measured independently of \(G\), the formula is not self-referential. The prediction \(\Omega_\Lambda = 2/3\) is testable without knowing \(G\): it requires only the dark energy fraction from expansion history (supernova luminosity distances, CMB flatness, baryon acoustic oscillations). Current measurements give \(\Omega_\Lambda \approx 0.685\), 2.7$\backslash$% above the prediction.

The prediction is not falsified. It is off by less than 3$\backslash$% using measurements that themselves carry systematic uncertainties of comparable magnitude. What the formula says is that a universe whose medium has only a vacuum floor and no matter would have \(\Omega_\Lambda = 2/3\) exactly; the presence of matter inflates the measured dark energy fraction slightly above this value.

\subsubsection{What This Does and Does Not Establish}
\label{sec:org9216477}

\emph{Established:} The structural form of \(G\) in Medium Theory terms. The coupling constant of the gravitational strain field is the medium's vacuum-floor compliance, \(H_0^2/\rho_\Lambda\), scaled by \(1/(4\pi)\). This is not a brute constant --- it is a ratio of medium properties. An equivalent testable prediction: \(\Omega_\Lambda = 2/3\). Agreement with two independent parameter sets at 94--97\\%.

\emph{Not yet established:} The full derivation of \(G\)'s SI value free of the Friedmann equation. The current measurement of \(\rho_\Lambda\) in \(\mathrm{kg\,m^{-3}}\) uses \(G\) as an input, making the formula partially self-referential at present. A clean derivation requires either an independent SI measurement of \(\rho_\Lambda\) or the formal strain field equations --- how the \(\iota\) packing geometry maps onto the inverse-square law, and how the projection boundary account gives \(G\) from substrate quantities without assuming the Friedmann equation. This is open work in the dependency chain, downstream of the strain field derivation.

\emph{What LQG cannot do:} Loop quantum gravity inherits \(G\). It has no mechanism for deriving it from within its own framework because \(G\) belongs to the level beneath the spin network. The Medium is pointing toward a derivation LQG cannot make.

\section{Unresolved Questions}
\label{sec:org1b47ca7}

These are the frontier of the framework, stated precisely and carried explicitly.

\textbf{The Combination Rules and the Algebraic Hierarchy.} A prior account derived a scalar hierarchy: \(\kappa = \sum \iota\) and \(\Kappa = \sum \kappa\), projecting to \(\hat{H} = \mathrm{proj}(\Kappa)\). That account assumed \(\iota\) produces scalar values combining by addition. The multifaceted \(\iota\) model — multiple relational modes, qualitatively distinct — makes scalar summation insufficient. You cannot sum qualitatively distinct things. The scalar hierarchy is a first approximation that captured one dimension of \(\iota\)'s structure. It requires rebuilding from the mode structure upward.

What the scalar hierarchy got right: the endpoint and the three-level structure. The Hamiltonian \(\hat{H}\) is the correct view-layer target. The claim that particle properties project from substrate combinatorics is correct. These survive revision.

What it got wrong: the combination operation. Addition is not the right rule for qualitatively distinct modes.

What the right rule is, I cannot say. The modes are different in kind, so combining them is more than adding numbers up. The rule has to be worked out, and it has to produce the particles we actually see.

The open question underneath this is whether the substrate --- binary characters plus the relational modes --- carries enough variety to reproduce the full Standard Model. I cannot settle that. It needs someone with deeper mathematical training than I have. I am marking it as the frontier, not claiming it is crossed.

What is not in question is the endpoint: whatever the rule turns out to be, it has to give back the Hamiltonian at the view layer.

\textbf{The Adjacency Constraint Under Revision.} The binary formulation — two \(\iota\) events of identical distinction character cannot be directly co-manifested --- is logically sound as stated. The open question is what it means when \(\iota\) events have internal mode structure. Two possibilities. First: identity means mode-identical across all relational modes, not just distinction character. Richer internal structure means fewer configurations are actually blocked. The constraint becomes less restrictive. Second: the constraint applies per mode — each mode has its own adjacency rule, and two \(\iota\) events adjacent in one mode but not another can still combine. The constraint becomes a family of rules. Either reformulation preserves the underlying logic — identical adjacent distinction produces nothing new --- while extending it to the multifaceted \(\iota\). The binary formulation is retained as the current working version. The mode-structure reformulation is the open derivation. Note: the adjacency constraint was carrying significant load in prior versions --- generating the particle table, explaining confinement, setting the minimum complexity threshold. All of that work is provisional pending the mode-structure reformulation.

\textbf{What the Field Measures.} The framework records distinction characters as \(\{+1,-1\}\). The original argument for binary: one act of distinction creates one boundary; one boundary has two sides. That argument assumed \(\iota\) is point-like --- making one clean cut, two sides, done.

The multifaceted \(\iota\) model undermines the assumption without necessarily undermining the result. A multifaceted \(\iota\) does not make a single cut. It makes a distinction with internal structure --- a boundary that has geometry, not just two sides. The question is no longer whether \(\iota\) is binary. It is what the field \(\phi\) actually measures.

The field may record the net output of a richer internal process. Like measuring voltage across a circuit: one number, one scalar, while the internal behavior is vastly more complex. The measurement collapses internal structure to a scalar output. \(\{+1,-1\}\) may be that scalar output --- the net distinction character --- while the relational modes are the internal mechanism the field does not directly see.

If that is right, the binary field is correct not because \(\iota\) is binary but because the field measures the result of distinction, not its internal structure. The modes govern interaction. The field records the outcome. These are separable claims.

Three possibilities remain open. First: the field is the full description --- \(\iota\) is genuinely binary and the mode structure is expressed entirely within \(\{+1,-1\}\). Second: the field is a net scalar projection of richer internal mode structure --- binary output, non-binary mechanism. Third: the field itself needs extension --- richer output values required once the mode structure is understood. Which of these is correct cannot be determined until the relational mode structure is derived. The binary field is retained as the working description. Its completeness is an open question.

\textbf{The Vocabulary Boundary.} The terms occurrence, registration, and boundary are proposed as genuinely primitive replacements for the inherited vocabulary of observer, measurement, and collapse. Whether they succeed --- whether they are themselves free of hidden Newtonian assumptions --- is an open question requiring formal examination.

\textbf{The Logic of Distinction.} Before a \(\iota\) event occurs, there is no fact about which character it will produce. The outcome is not unknown --- it does not yet exist. A distinction becomes definite only in the act, not before it. This is a claim about when an outcome becomes determinate, not about how many aspects \(\iota\) has. Whether it matches any existing formal logic is an open question. I suspect it does. I cannot prove it.

\textbf{The Nature of \(\iota\).} What is intrinsic distinction? The provisional characterization --- the property permitting local, spontaneous reduction from maximum undifferentiation in a fully random substrate --- is the central open question of the theory.

\textbf{The Formal Account of the Projection Boundary.} How does an observer embedded in the binary field construct a dimensional coordinate system? What determines which dimensional categories emerge? This is the primary task for the next stage of formal development.

\textbf{The Projection Boundary Derivation --- A Dependency Chain.} The framework claims one \(\iota\) event = 1 in Medium units and that \(\hbar\) is the SI measurement of that unit. That claim is significant: it asserts that the quantum of action corresponds to one act of substrate self-distinction. Until it is derived, the correspondence is an empirical match, not a result.

The derivation requires showing that an embedded observer constructing a dimensional coordinate system from sustained \(\iota\) configurations will necessarily assign the action dimension --- energy multiplied by time --- to a single \(\iota\) event. Why that dimensional combination and not another is the core question. It cannot be answered until the observer's dimensional categories are derived from the relational web structure. That derivation requires knowing how configurations project to dimensional quantities. Which requires the formal projection boundary account. Which requires the combination rules. Which requires the \(\iota\) mode structure.

The projection boundary derivation is downstream of everything else. It is not missing because it has been neglected. It is missing because it sits at the end of a logical dependency chain:

\(\iota\) mode structure → combination rules → stable configurations → projection boundary account → action dimension derivation → \(\hbar\) as derived result

Each step in the chain is a named open derivation. The chain is the research program. Completing it in sequence is the path from empirical match to derived result.

Three further derivations depend on the same chain. Why projection fixes the ratio of projected distance to projected time at the specific value observers measure as \(c\) --- \(c\) is view-layer, not substrate, but why projection fixes it at that value requires the projection boundary account. Why the minimum projected mass unit takes its specific value requires the gravity account, which requires the strain field derivation, which requires the mode structure.

\(G\) has a partial derivation: \(G = H_0^2/(4\pi\rho_\Lambda)\), giving 94--97\\% agreement and the equivalent prediction \(\Omega_\Lambda = 2/3\) (see Derivation III). What remains is the independent measurement of \(\rho_\Lambda\) in SI units --- currently obtained via the Friedmann equation with \(G\) as an input --- and the formal strain field equations connecting \(\iota\) packing geometry to the inverse-square law. Both sit downstream of the projection boundary account.

The substrate has no preferred frame. Frames arise only at projection. Lorentz invariance is a property of how projection works. The derivation confirming that projection recovers exact Lorentz invariance sits at the end of the same chain.

These are not independent gaps. They are one gap with a known structure. The dependency chain is the shape of the open work.

\textbf{The Observation Method Problem.} The Standard Model's particle table was built primarily from high-energy collision experiments. These experiments force the view layer into configurations above the stability threshold --- energy conditions that do not occur in stable physical reality. The particles produced are almost entirely transient, decaying within fractions of a second into stable configurations. The Standard Model catalogs the full debris of those collisions as fundamental particles. The Medium does not accept that catalog without examination.

The framework distinguishes two classes of measurement. The first: naturally stable configurations --- the electron, proton, photon, neutrino, and the lightest stable hadrons. These are candidates for genuine stability extrema in the substrate. They persist. They are not products of forced above-threshold conditions. The second: transient collision products --- particles that exist only under extreme forced conditions and decay immediately. These are candidates for above-threshold view-layer forced states, not stable substrate configurations. They reveal the path to stability; they are not the stable configurations themselves.

This distinction matters for the derivation program. The Medium is not required to derive every particle the Standard Model catalogues from collider debris. It is required to derive the naturally stable configurations --- and then to show why forced above-threshold conditions produce the transient states observed. The Standard Model treats both classes equally as fundamental. The Medium should not.

Neutrino oscillation sits partially in the first class. Solar, atmospheric, and reactor neutrinos are natural. Their oscillation is observed over long baselines, not in collision debris. That makes neutrino oscillation a more reliable guide to substrate structure than most collider-produced phenomena. The PMNS mixing parameters, however, are partly measured in accelerator experiments, and whether those values are substrate properties or measurement-method artifacts remains open.

\textbf{Particle Configurations.} Stable configurations are those where the relational web reaches a stability extremum under the combination rules governing the \(\iota\) relational modes. The open derivation: identify the minimum stable \(\iota\) patterns, match them to the naturally stable observed particles, and show why forced above-threshold conditions produce the transient states the Standard Model catalogs as additional particles. Note: the adjacency constraint and scalar hierarchy that previously defined this derivation are under revision pending the mode structure derivation. The particle table derivation is provisional until the combination rules are established. The direction is current; the formal path is open.

\textbf{The Higgs Mechanism.} The Standard Model gives particles mass through coupling to the Higgs field. Coupling strength determines mass. The Higgs boson was discovered --- confirmed physics, a measured particle. The Medium says mass corresponds to \(N\) as a first approximation and says nothing about the Higgs. That silence is a gap. Three possible placements. First: the Higgs is view-layer --- the observer's account of why different substrate configurations have different \(N\) counts, with the Higgs field being the view-layer description of that substrate structure. This would make the Higgs a projection artifact of the \(N\) structure, but it would need to explain why the Higgs has its specific mass and why coupling strengths match the observed fermion mass spectrum. Second: the Higgs is substrate-emergent --- a specific \(\iota\) configuration that permeates the Medium and whose co-manifestation with other configurations determines their effective \(N\). This preserves the Higgs as a real phenomenon at the substrate level but introduces a new substrate entity not yet in the framework. Third: the placement is genuinely open. The Higgs mechanism is confirmed physics. The Medium does not yet address it. This is the current position.

\textbf{The Quantum Vacuum.} Sub-threshold substrate combinations appear at the view layer as quantum vacuum fluctuations. The quantum vacuum is not empty --- it exerts measurable pressure and perturbs energy levels. The Casimir effect: two conducting plates in vacuum experience an attractive force from vacuum fluctuations, quantitatively confirmed. The Lamb shift: vacuum fluctuations shift hydrogen energy levels by a precisely measured amount, a result that drove the development of QED. Both are view-layer QED results. Whether the substrate fluctuation account produces the correct vacuum pressure and energy perturbation for these measurements is not shown. The account is directionally consistent --- sub-threshold combinations are real substrate structure --- but the quantitative bridge from substrate to measured Casimir force and Lamb shift does not yet exist.

\textbf{Adjacency Constraint and Pauli Exclusion.} The adjacency constraint operates at the substrate level: two \(\iota\) events of identical distinction character cannot be directly co-manifested. Pauli exclusion operates at the view layer: no two fermions can occupy the same quantum state. They are at different levels and are not the same rule. The relationship between them is unexamined. The possibility worth pursuing: the adjacency constraint may generate Pauli exclusion under projection. If identical distinction characters cannot be directly co-manifested at the substrate, configurations that project as fermions --- half-integer relational closure --- may inherit that exclusion at the view layer. That would be a derivation of Pauli exclusion from substrate structure. The gap: the adjacency constraint applies to all \(\iota\) events regardless of what they project to. Bosons do not have Pauli exclusion. Why the constraint produces exclusion for fermions but not bosons is not shown. Until that is shown, the adjacency constraint and Pauli exclusion are parallel structures at different levels, not the same result at different scales. This is a named candidate derivation, carried as open.

\textbf{Spin as Relational Closure.} Which encounter patterns produce which spin values (0, 1/2, 1, 3/2, 2)? The claim that spin is a property of relational closure --- integer spin closing in one traversal of the web, half-integer in two --- is current. The derivation of which specific encounter patterns correspond to which values is open. More consequentially: the framework names the topology of spin but does not yet connect it to exchange statistics. Fermions obey Pauli exclusion --- no two in the same quantum state. Bosons do not. That difference determines atomic shell structure, the stability of matter, and the behavior of light. The spin-statistics theorem in quantum field theory derives this connection from Lorentz invariance, causality, and positivity of energy. The Medium's relational closure description is consistent with the topology the theorem requires but does not derive the statistics. Doing so requires showing that projection from the substrate recovers Lorentz invariance, from which the spin-statistics theorem follows as a downstream result. That derivation does not yet exist.

\textbf{Derivation of Prior Frameworks.} Three existing frameworks are candidates for derivation from the substrate. Planck bits: the registered encounter character is a bit --- the derivation shows where the bit comes from. Loop quantum gravity: the encounter web maps to spin networks; stable configurations map to nodes; encounter relationships map to edges; the derivation shows LQG as the observer's projection of the substrate's relational structure. String theory: different encounter configurations at stability extrema correspond to different particles as string vibration modes correspond to different particles --- the derivation requires showing how the encounter structure generates the dimensional space string theory presupposes. Each is an open derivation. Each is a named target.

\textbf{Fractional Charge.} The framework produces integer net distinction characters from ±1 combinations. Quarks carry charge ±1/3 and ±2/3 --- never observed in isolation, always confined within hadrons. The hadronic totals are integers; the quark-level values are not.

The appearance of fractional charge in a binary substrate is not a puzzle to be solved by making the fractions work. It is a signal --- the same kind of signal dark energy provided. Dark energy appeared to require explanation until the framework asked why it appears at all, found it was a projection artifact of misread redshift, and dissolved it. The same instinct applies here.

A binary substrate produces integers. The appearance of 1/3 means the quark model is describing substrate structure at the wrong level --- carving the substrate at the wrong joints, the same way recession velocity was the wrong description of redshift. Quarks as described in the standard model may not be the right substrate entities. The framework should ask what configuration of \(\iota\) events produces the behavior attributed to quarks. That configuration will have integer distinction characters. The fractional values appear because the standard model imposes its own framework on substrate structure that has no fractional properties.

The open derivation: identify the \(\iota\) configurations whose behavior the standard model describes as quarks. Show that those configurations carry integer characters. Show why projecting through the standard model's quark framework produces the appearance of 1/3 and 2/3. Confinement --- the impossibility of isolating a quark --- is then not a separate physical fact but a structural consequence of the same misidentification: outside the hadronic configuration, the entity the standard model calls a quark does not exist as a substrate unit.

Treating fractional charge as something that exists only inside a proton or neutron --- not as a free-standing property --- is a candidate direction for understanding why the fractions appear and sum correctly. It is not the primary resolution. The primary resolution is the correct substrate-level identification of the configurations the standard model is describing.

\textbf{The Mass Hierarchy.} Mass corresponds to \(N\) as a first approximation. That correspondence does not explain why the muon and tau are far heavier than the electron, though identical to it in structure, charge, and spin. It does not explain why there are exactly three of them. The Standard Model inserts these values by hand. The Medium claims mass is derivable from substrate structure --- that claim requires showing which combinatorial configurations produce which specific \(N\) counts, and why those counts correspond to the observed values. This derivation does not yet exist.

\textbf{The Cosmic Microwave Background and Big Bang Nucleosynthesis.} Rejecting the Big Bang interpretive framework creates a precise obligation: provide alternative accounts of the observations that framework explains. The observations are real. They are not in dispute. The interpretive claims built on them are what the foundation challenges. That distinction requires a clear accounting of what is and is not yet addressed.

The CMB presents three distinct constraints. They are not equal.

\emph{Uniformity:} The CMB temperature is uniform across the sky to one part in 100,000. The horizon problem section addresses this directly: the substrate has no preferred location, \(\iota\) events are statistically uniform across the relational web, and that homogeneity projects as uniform temperature. No causal contact required. This constraint has a candidate dissolution.

\emph{The blackbody spectrum:} The CMB follows a near-perfect blackbody curve at 2.725 K. A blackbody spectrum requires thermal equilibrium. The standard account achieves this through a hot opaque early phase. The Medium has no hot early phase. An alternative account requires showing that the iota medium reaches or maintains thermal equilibrium at this temperature through a different mechanism. No such account exists. This is an open obligation.

\emph{The anisotropy spectrum:} The one-in-100,000 fluctuation pattern matches predictions of acoustic oscillations in an early baryon-photon plasma with extraordinary precision. This is the strongest CMB constraint. It is not addressed. It is not dissolved. It requires a substrate account of which \(\iota\) configurations produce which fluctuation amplitudes at which angular scales. That derivation does not exist. This is the primary open obligation in the cosmological section.

Big Bang nucleosynthesis is a separate and independent constraint. The observed amounts of the light elements --- hydrogen, helium, and traces of others --- are reproduced by applying known nuclear physics to the conditions of the early universe. The precision is extraordinary. This is not interpretation layered on observation --- it is nuclear physics producing a specific numerical prediction that matches measurement. The Medium offers no alternative account of these abundances. This is the single strongest observational constraint on any alternative cosmology. It cannot be dissolved by challenging interpretive frameworks. It requires a physical account of how those specific ratios arise. That account does not exist.

The accounting is precise: uniformity has a candidate dissolution; the blackbody spectrum and anisotropy spectrum are open obligations; BBN abundances are the hardest unsolved problem in the framework. All three are carried forward. None are dismissed.

\textbf{Color Charge --- Dissolved.} Color charge is never directly observed. Confinement guarantees that no colored object reaches a detector. Every observed state is color-neutral. The label was borrowed from optics for one reason: red + green + blue = white. Three things combine to produce neutrality. That is the entire justification for "color." The physics is the neutrality --- the three-fold combination rule that produces stability. The optical analogy adds nothing. It names the pattern without explaining it.

The SU(3) gauge freedom confirms dissolution from the Standard Model's own structure: the freedom to relabel red/green/blue without changing any observable is the mathematical statement that color is not a physical property. It is a labeling convention.

Gauge symmetry survives and gains a substrate ground. The iota's relational modes have no preferred global labeling --- the medium has no preferred orientation of modes. That is the same freedom SU(3) describes, now located at the substrate rather than imposed as an external requirement on a field. Color charge dissolves. The gauge freedom that generated SU(3) does not. It emerges from the iota mode structure. Whether the iota mode structure generates exactly SU(3) or something more general from which SU(3) is a projection is the open derivation.

The iota relational mode model makes color charge redundant at both levels. If the iota has multiple relational modes combining under rules that produce stable two-body and three-body configurations, the stability pattern falls out of the mode structure directly. Color charge is not a substrate property requiring placement. It is a view-layer label for a substrate pattern the iota model accounts for without it. Color charge dissolves as a Standard Model artifact of describing hadronic stability from collider observations rather than from substrate structure.

\textbf{Dark Matter.} Dark matter interacts gravitationally but not electromagnetically, not via the strong force, and at most weakly. It is invisible, stable, and accounts for roughly five times more mass than visible matter. In The Medium's terms, a dark matter candidate is a stable \(\iota\) configuration at a stability extremum with nonzero \(N\) (mass --- gravitational coupling at the view layer), zero net distinction character (electrically neutral), and no relational closure that engages the projection channels producing electromagnetic or strong force coupling. It is a pattern that achieves stability under the combination rules but does not light up the observer's EM or strong projection channels. From the substrate perspective, dark matter is not mysterious --- it is a class of stable configurations that happen not to couple to the forces visible matter couples to. The open derivation: identify which specific \(\iota\) patterns satisfy the combination rules, achieve a stability extremum, carry zero net distinction character, and have no EM or strong relational closure. Those patterns are the dark matter candidates. Whether the framework produces candidates matching the observed mass distribution and halo structure is the empirical test.

\textbf{Color Charge, Gravity, and the Unassigned Category.} Color charge, gravity, neutrino oscillation, and dark matter have not been placed with confidence in either layer. Gravity is the most consequential.

General relativity identifies gravity with spacetime curvature. The Medium makes spacetime view-layer emergent. If spacetime is view-layer, gravity as spacetime curvature is also view-layer. That placement has three consequences worth naming.

First, a potential strength: the equivalence principle --- inertial mass equals gravitational mass to extraordinary precision --- may be automatic. If mass corresponds to \(N\) and \(N\) governs both inertial resistance and gravitational coupling, the equivalence principle is a structural consequence of one substrate quantity projecting two ways. It stops being a mysterious coincidence.

Second, an open problem: quantum gravity. Every other force has been quantized. Gravity resists. If all four forces are view-layer, quantum gravity remains a view-layer problem and The Medium does not obviously resolve it. The possible path: gravity may not be like the other forces. It may bridge both layers --- substrate configuration density shaping projection geometry, projection geometry feeding back on substrate structure. That bridge is the candidate direction. It is not developed.

Third, gravitational waves propagate through spacetime. If spacetime is view-layer, gravitational waves are view-layer phenomena. LIGO detects them. That is consistent with the placement.

Whether color charge is substrate or view-layer emergent; how dark matter candidates arise from substrate structure --- these remain at the frontier. Dark energy is not a frontier question. It dissolves when redshift is correctly understood.

A candidate direction for gravity emerged from the iota packing model (developed 2026-05-23). Space is relational --- exists as relationships between \(\iota\) events, not as a void or independent fabric. Spacetime curvature is a large-scale pattern of \(\iota\) alignment. A gravity well is a persistent local deviation in packing that the surrounding medium tries to correct.

Mass is the stable stochastic imperfection. Gravity is the strain field. The \(\iota\) packing model gives infinite reach without a carrier particle: every \(\iota\) event participates in the ground state, so every stable imperfection propagates strain through the medium without limit and without signal. The equivalence principle is automatic --- the same \(N\) that resists inertial change is the \(N\) whose packing imperfection generates the strain field.

Quantum gravity's resistance to derivation is consistent with this account. Gravity is not like the other forces at the view layer because it is not a view-layer force in the same sense. It bridges substrate packing geometry and view-layer spacetime. Attempting to quantize it applies view-layer tools to something that only exists because view-layer granularity was averaged away. The failure is structural, not mathematical.

This direction requires: showing how the strain field through the \(\iota\) medium maps onto the inverse square law; determining what constrains the iota's relational mode count; and connecting the packing geometry to the projection mechanism that produces the spacetime manifold. None of these derivations exist. The direction is named.

\textbf{The Observer Boundary.} The observer is a sustained binary configuration within \(\phi\), attempting to describe the substrate that produced it from within. Every formal system contains truths it cannot prove from within itself. The Medium faces an analogous boundary and acknowledges it.


\section{Summary}
\label{sec:org1053bb3}

\begin{table}[htbp]
\caption{Complete symbol table for version 0.5}
\centering
\begin{tabular}{llp{4.5cm}l}
\toprule
\textbf{Symbol / Concept} & \textbf{Definition} & \textbf{Notes} & \textbf{Status}\\[0pt]
\midrule
\(\iota\) & The Medium's act of self-distinction --- produces character as the difference the distinction makes & Not point-like; multiple relational modes; intrinsic cohesive attraction & Axiomatic\\[0pt]
\(\phi : M \to \{+1,-1\}\) & Distinction character field --- records net output of \(\iota\) events & May be scalar projection of richer mode structure; completeness open & Current\\[0pt]
\(x \sim y\) & Co-manifestation --- shared state, no spatial relationship & Prior to locality & Current\\[0pt]
Relational modes & Qualitatively distinct aspects of the \(\iota\) event governing interaction & Count undetermined; constrained by what framework must produce & Open\\[0pt]
Combination rules & Rules governing how \(\iota\) relational modes combine & Expected form: constrained group-theoretic algebra on tensor product of mode states & Open\\[0pt]
\(\hat{H} = \mathrm{proj}(\cdot)\) & Hamiltonian --- projection of stable relational web configurations into spacetime & Endpoint of combination rules derivation; derived not assumed & Current\\[0pt]
\(\kappa = \sum \iota\) & Scalar coupling character --- net distinction character of a local cluster & First approximation; provisional pending mode combination rules & Provisional\\[0pt]
\(\Kappa = \sum \kappa\) & Scalar configuration energy --- governs stability via extremum condition & First approximation; provisional pending mode combination rules & Provisional\\[0pt]
Mass & Stable stochastic imperfection in the \(\iota\) packing & \(N\) (iota count) is view-layer first approximation only & Current\\[0pt]
Gravity & Strain field through the \(\iota\) medium toward ground-state restoration & Infinite reach without carrier particle; candidate direction & Current\\[0pt]
\(\mathrm{sgn}(\Kappa)\) & Net distinction character of configuration --- projects to charge & Charge neutrality = exact character balance & Provisional\\[0pt]
Spin & Relational closure of configuration's encounter pattern within the web & Integer: one traversal; half-integer: two; derivation open & Provisional\\[0pt]
Adjacency constraint & Identical adjacent distinction characters forbidden & Binary formulation current; mode-structure reformulation open & Provisional\\[0pt]
Projection boundary & Planck scale --- where pre-dimensional substrate meets dimensional observer & Correspondence with \(\hbar\) is empirical match pending derivation & Current\\[0pt]
\(\iota_{\mathrm{proj}} = 1\) & One projected \(\iota\) event --- the Medium's unit of action & \(\hbar\) is the SI measurement; action dimension requires derivation & Current\\[0pt]
\(c\) & Projection-layer ratio of projected distance to projected time & Not substrate-level; why projection fixes this value is open & Provisional\\[0pt]
\(G\) & Gravitational constant --- vacuum-floor compliance of the strain field: \(G = H_0^2/(4\pi\rho_\Lambda)\) & 94--97\\% agreement; equivalent prediction \(\Omega_\Lambda=2/3\); full derivation pending formal strain field equations & Partial\\[0pt]
\(x \sim y \to\) entanglement & Co-manifestation projects to quantum entanglement & Mystery dissolves at substrate level & Current\\[0pt]
Wave function & View-layer representation of relational web probability landscape & Born rule is how view layer reads probability density & Current\\[0pt]
Wave function collapse & View layer resolving one reading from probability landscape & Not a physical event & Dissolved\\[0pt]
Entropy & Projection-layer emergent --- observer's coarse-graining of microstates & No substrate analog & Current\\[0pt]
Time & Projection-layer emergent --- ordered account of substrate transitions & No substrate universal clock & Current\\[0pt]
Redshift & Three-layer projection artifact --- configuration-depth asymmetry & Hubble linearity requires substrate gradient derivation & Current\\[0pt]
Color charge & Labeling convention for three-fold combination stability pattern & Dissolves; SU(3) gauge freedom survives in iota mode structure & Dissolved\\[0pt]
Dark matter & Stable \(\iota\) configuration: nonzero mass, no EM/strong relational closure & Candidate sketch; derivation open & Open\\[0pt]
Dark energy & Projection artifact of misread redshift & Not a substrate phenomenon & Dissolved\\[0pt]
GR/QM conflict & Category error --- two frameworks describing different layers & GR: view-layer smooth aggregate; QM: substrate granularity at view layer & Dissolved\\[0pt]
\bottomrule
\end{tabular}
\end{table}


\section{A Personal Note}
\label{sec:orgf229b25}

I am seventy-nine years old. I am not an academic. I have no credentials in physics or mathematics. I began this work with a question that would not leave me alone and a conviction that the question deserved a serious attempt at an answer.

I began this knowing I may not live to see it reach its conclusions. That does not trouble me. What matters is that the ideas are documented, attributed, and made available to anyone who finds them worthy of development. If this framework contains something true it belongs to whoever can take it further. My name attached to it is enough --- for the record, and for my family.

Gregor Mendel published his work on inheritance and it sat unread for decades before being recognized as foundational. The ideas belong to whoever finds them. The record belongs to their origin.


\section{License and Attribution}
\label{sec:org4b9078a}

This work is licensed under the Creative Commons Attribution 4.0 International License (CC BY 4.0). You are free to share and adapt this material for any purpose, including commercially, provided you give appropriate credit to the original author, provide a link to the license, and indicate if changes were made.

\url{https://creativecommons.org/licenses/by/4.0/}

All versions of this document are archived on Zenodo and remain linked to the authorship record of Bradley K. DePuy, independent researcher.

DOI: 10.5281/zenodo.18927154\\[0pt]
ORCID: \url{https://orcid.org/0009-0001-0328-2299}
\end{document}
